% 本文件由 LaTeX 排版 Agent 根据有效 translation.json 自动生成。
% author=Councils; work_id=3816; title=迦太基会议（公元419年）

\chapter{迦太基会议（公元419年）}

% structure: p0001 source=h1 target=omitted reason=page-title-already-rendered-by-parent

\Emph{普遍称为\WorkTitle{非洲教会法典}}。\par

% structure: p0003 source=h2 target=section reason=relative-heading-rank; base=h2

\section{引言注释}

\Quote{在初期时代，没有哪个地方比非洲召开会议更频繁。在418-19年，以前在迦太基举行的十六次会议、一次在米莱维、一次在希波所制定的所有教规，凡获批准者，均被宣读，并得到当时在迦太基召开会议的众多主教们的新认可。这部汇编就是\WorkTitle{非洲教会法典}，它在所有教会中一直享有仅次于\WorkTitle{普世教会法典}的最高声誉。这部法典在古代英格兰教会中具有极大权威，因为埃格伯特的许多\WorkTitle{埃格伯特摘录}都转录自它。尽管\WorkTitle{普世教会法典}以迦克墩教规结束，但这些非洲教规被插入东西方教会的古代法典中。这些教规虽经一次主教会议的批准和认可，但似乎被某个无知的私人手笔划分或编号，并且可能遇到了非常不熟练的抄写员，因此其中一些教规的意义和连贯性变得混乱不清。狄奥尼修斯·爱克西古斯等人将它们称为\Emph{\WorkTitle{非洲主教会议教规}}。虽然最初并非同时制定，但它们都得到了一次非洲主教会议的确认，这些主教在诵读了信经和尼西亚大公会议的二十条教规之后，继续制定新教规，并强化旧教规。}\par

迦太基原是整个非洲的首府，正如圣奥斯定在他的\WorkTitle{书信}第162号中所告诉我们的。正因如此，在那里从非洲各省聚集召开了大量会议。特别是在大主教奥勒留坐镇期间，这些主教们的集会频繁举行；通过这些会议，为在非洲建立教会纪律制定了许多教规。最后，在奥古斯都贺诺留（第十二任执政官）和狄奥多西（第八任执政官）执政之后，于六月朔日前第八日，即我们的主后419年5月25日，在同一城市召开了另一次会议，会上审议了先前通过的所有教规，其中大部分再次得到了主教会议的权威确认。这些经此会议确认的教规，从那天起至今被称为\Quote{\WorkTitle{非洲教会法典}}。这些教规最初并非以希腊文而是以拉丁文通过，并以同一语言确认。狄奥尼修斯·爱克西古斯在其序言中明确见证这一点，该序言附于包含这些教规的\Quote{\WorkTitle{教会教规法典}}中。不确定这些迦太基会议的教规何时被译成希腊文。唯一可以确定的是，它们在楚洛会议之前已被译成希腊文，楚洛会议在其第二条教规中将这些教规收入希腊\WorkTitle{教规汇编}，并得到了该会议的权威确认；从此以后，这些教规在东方教会中与其他一切教规享有同等地位。\par

有一个极其有趣的问题：这部汇编作为一部汇编的权威性是什么，以及这部汇编是如何形成的？似乎毫无疑问，我们所知的这部汇编基本上是楚洛会议所接受的法规，其后第七次大公会议（第二次尼西亚会议）给予的普遍认可赋予了这些教规准大公会议的权威。\par

以下是制定这些法规的各次会议及其日期。\par

\begin{itemize}
\item 迦太基（格拉图斯主持）——公元345-348年
\end{itemize}

\begin{itemize}
\item （格内特利乌斯主持）——387或390年
\end{itemize}

\begin{itemize}
\item 希波——393年
\end{itemize}

\begin{itemize}
\item 迦太基——394年
\end{itemize}

\begin{itemize}
\item （6月26日）——397年
\end{itemize}

\begin{itemize}
\item （8月28日）——397年
\end{itemize}

\begin{itemize}
\item （4月27日）——399年
\end{itemize}

\begin{itemize}
\item （6月15日）——401年
\end{itemize}

\begin{itemize}
\item （9月13日）——401年
\end{itemize}

\begin{itemize}
\item 米莱维（8月27日）——402年
\end{itemize}

\begin{itemize}
\item 迦太基（8月25日）——403年
\end{itemize}

\begin{itemize}
\item （6月）——404年
\end{itemize}

\begin{itemize}
\item （8月25日）——405年
\end{itemize}

\begin{itemize}
\item （6月13日）——407年
\end{itemize}

\begin{itemize}
\item 迦太基（6月16日和10月13日）——408年
\end{itemize}

\begin{itemize}
\item 迦太基（6月15日）——409年
\end{itemize}

\begin{itemize}
\item （6月14日）——410年
\end{itemize}

\begin{itemize}
\item （5月1日）——418年
\end{itemize}

\begin{itemize}
\item （5月25日）通过了\WorkTitle{非洲法典}——419年
\end{itemize}

非洲会议的编号在不同作者之间差异很大，凯夫认为在401至608年间有九次，在215至533年间有三十五次迦太基会议。在布伦斯的\Emph{\WorkTitle{宗徒及古代大公会议教规选辑}}结尾处可以找到非常实用的表格，显示了富克斯的结论。\par

% structure: p0028 source=h2 target=section reason=relative-heading-rank; base=h2

\section{古代引言}

（见于狄奥尼修斯·爱克西古斯\WorkTitle{教规法典}，米涅\WorkTitle{拉丁教父文集}第67卷，第182栏。）\par

在最荣耀的皇帝奥古斯都贺诺留（第十二任）和狄奥多西（第八任）执政之后，六月朔日前第八日，在迦太基的福斯泰大殿的议事厅内，当主教奥勒留就座，与努米底亚省的首席主教瓦伦提尼、意大利皮切诺省的波腾提内教堂的福斯蒂诺（罗马教会特使）、以及来自不同非洲各省的特使，即两个努米底亚省、拜扎凯纳省、凯撒利亚毛里塔尼亚省和特里波利省的代表，还有文森特·科洛西塔努斯、福图纳提安及行省总督区的其他主教，共二百一十七人，以及罗马教会的司铎兼特使腓力和阿塞鲁斯，执事侍立一旁，主教奥勒留说等等，\Emph{如下文}。\par

% structure: p0031 source=h2 target=section reason=relative-heading-rank; base=h2

\section{教规}

主教奥勒留说：你们最蒙福的弟兄们，请记得，在指定召开主教会议的日子之后，我们等候我们那些已被派为代表并已抵达本次会议的弟兄们时，讨论了许多事情，这些都必须记录在案。因此，让我们因如此众多会众的聚集而感谢我们的主。现在应当提出我们已有的尼西亚主教会议的法案，这些由教父们所决定，以及我们在此的前辈们所制定的那些事项，他们确认了同一次会议，或者按照同样形式，由各级圣职人员（从最高直至最低）有益地制定的事项。全体主教会议说：请提出它们。\par

公证员达尼尔宣读：尼西亚主教会议的信德宣认或法令如下。\par

当他正在说话时，\Person{Faustinus}，\Place{Potentia}地方的人民的主教、意大利\Place{Picenum}省人，罗马教会的特使说：\Quote{宗座已托付给我们若干书面文件，另有一些如同要跟你们的真福讨论的训令，正如我们在上述记录中所追忆的，即关于在\Place{尼西亚}所订的教规，要遵守其法令和习俗；因为有些是出于法令和教规而遵守，有些则是出于习俗。因此，首先就这些事，若蒙你们的真福同意，让我们查问；然后，其他已采纳或拟议中的训令才可被确认；这样你们就能通过回函向宗座表明，并向同一位可敬的教宗宣告，我们已经殷勤地记住了这些事；尽管行动标题已载入记录中。在这事上，我们应当如我上面所说，按你们所爱的真福所愿而行。因此，让备忘录呈到中间，使你们能识别其中所包含的内容，以便可以对每一点给予答复。}\par

\Person{Aurelius}说：\Quote{让备忘录被提出，这是我们的弟兄和同僚最近放入记录中的，然后其余已做或要做之事依次进行。}\par

公证人\Person{Daniel}宣读了备忘录：\Quote{致我们的弟兄\Person{Faustinus}和我们的儿子、司铎\Person{Philip}与\Person{Asellus}，主教\Person{Zosimus}。你们清楚记得我们曾将某些事务托付给你们，现在（我们吩咐你们）要执行一切，如同我们自己在那里一样，因为，确实，我们的临在是与你们同在的；尤其你们有此命令，以及为了更大确定而插入本备忘录中的教规言辞。因为我们的弟兄们在\Place{尼西亚}大公会议中就这样说，当他们制定这些关于主教上诉的法令时：}\par

\Quote{但若一位主教曾被控告，等等。} \EditorialNote{此处逐字逐句引用\OriginalTerm{撒尔底迦}{Sardica}第五条教规。}\par

\Emph{古代摘要：若主教们已罢免了一位主教，而他向罗马主教上诉，则应被仁慈地听取，罗马主教写信或下令。}\par

当此被宣读后，\Place{塔加斯特}教会的主教、\Place{努米底亚}省的特使\Person{Alypius}说：\Quote{关于此事，在我们议会的前几次会议中已有一些立法，我们声称将永远遵守\Place{尼西亚}大公会议所规定的内容；然而我记得，当我们查阅此\Place{尼西亚}会议的希腊抄本时，并未找到所引用的这些话——为何如此，我确实不知。因此我们恳请您的尊敬、圣主教\Person{Aurelius}，鉴于据说\Place{尼西亚}大公会议法令的可靠文本保存于\Place{君士坦丁堡}城，请您惠允派遣使者，带着您圣座的信函，不仅致我们至圣的弟兄\Place{君士坦丁堡}主教，也致尊敬的\Place{亚历山大}和\Place{安提约基雅}主教，他们将把该大公会议的法令连同其签名的认证寄给我们，以便今后一切歧义得以消除，因为我们未能找到我们的弟兄\Person{Faustinus}所引用的言辞；尽管如此，我们承诺短期内遵循这些言辞，如我所说，直到可靠的抄本到手。此外，也应请求\Place{罗马}教会的尊敬主教\Person{Boniface}同样惠允派遣使者到上述教会，他们应根据其回函持有相同的抄本，但我们将我们手上有的上述\Place{尼西亚}大公会议抄本放入这些记录中。}\par

\Person{Faustinus}主教、\Place{罗马}教会的特使说：\Quote{愿您的圣洁不要在此事或任何其他事上羞辱\Place{罗马}教会，怀疑教规，如同我们的弟兄和同僚主教\Person{Alypius}所提及的；但请您惠允将这些事写给我们圣洁和至福的教宗，以便他寻出真正的教规，就所规定的一切与您的圣洁商议。但足够的是，最蒙福的\Place{罗马}城主教应进行查问，如同您的圣洁方面所提议的，以免教会之间似乎产生了任何争论，反而使你们更有能力在听到他的意见后，以兄弟般的爱德商议应遵守的最佳方案。}\par

\Person{Aurelius}主教说：\Quote{除了记录中已载明的，我们还将以我们卑微者的信函，更充分地告知我们圣洁的弟兄和同僚主教\Person{Boniface}我们所考虑的一切。因此，若我们的计划合大家之意，请让我们从众人口中得知。}全体议会说：\Quote{此议合我们之意。}\par

\Person{Novatus}主教、\Place{毛里塔尼亚Sitifensis}的特使说：\Quote{我们现在想起此备忘录取有一项关于司铎和执事的内容，即他们应由自己的主教或邻近主教审判，而这我们在\Place{尼西亚}大公会议中并未找到。因此，请您的圣洁命令此部分被宣读。}\par

\Person{Aurelius}主教说：\Quote{让所要求之处被宣读。}公证人\Person{Daniel}宣读如下：\Quote{关于圣职人员的上诉，即较低级别者，此同一届议会有一个确定的答复，关于此事你们应如何行，我们认为应插入如下：}\par

\Quote{\Person{Hosius}主教说：我不应隐瞒直到如今进入我心中的想法。若某位主教偶有快速动怒（这是不应发生的事），并迅速或严厉地对待其司铎或执事，并想要将他逐出教会，则应作出规定，使无辜者不被定罪，或被剥夺共融；被逐者应有权向相邻教区的主教上诉，使他的案件得到审理，并且因为不应拒绝请求听审的人，应更加勤勉地进行。而那位或公正或不公正地拒绝他的主教，应耐心容许此事被讨论，以便他的判决或被批准，或被修改，等等。}\par

这篇文件宣读后，努米底亚省希波教会主教奥古斯丁说：我们承诺将遵守此文件，前提是在更仔细的审查后发现它确实来自尼西亚会议。主教奥勒留说：如果这也令各位的爱德满意，请以你们的投票予以确认。全体会议说：尼西亚会议所规定的一切都令我们满意。苏菲图拉教会主教、拜扎塞纳省代表约昆杜斯说：尼西亚会议所制定的法令在任何细节上都不能被违背。\par

罗马教会代表福斯蒂努斯主教说：就阁下及圣阿利皮乌斯、我们弟兄约昆杜斯的宣认所能看出，我认为某些要点被削弱了，另一些则被确认了，这是不应发生的，因为连这些法令本身也受到了质疑。因此，为使我们在您和各位之间和谐一致，务请阁下屈尊将此事提交给神圣可敬的罗马教会主教，使他能考虑是否应准许圣奥古斯丁所颁布的内容，我是指关于下级上诉事宜。因此，如果在此事上仍有疑问，应将最蒙福之座的主教告知，如果这一点能在已获批准的法令中找到的话。\par

主教奥勒留说：正如我们已向各位爱德所建议的，请允许将尼西亚会议法令的副本以及那些由我们的前任们按照该会议的规定在这座城市中最有益地制定的事宜，以及我们如今所规定的事宜，一并宣读并载入会议记录。全体会议说：信经的副本、尼西亚会议的法令（这些文件以前是由你们圣座那位蒙福的、曾与会的前任采奇利安带到我们会议的），以及那些教父们随后在这座城市中制定的法令副本，或是我们现在经过共同商议所制定的法令，应被载入这些教会会议记录中，以便（如前所述）你们的圣座惠然致函最可敬的安提约基雅、亚历山大里亚和君士坦丁堡教会，请他们寄出最准确的尼西亚会议法令副本，并附上他们的签名认证；如此，事实真相得以澄清后，我们的弟兄、现任同僚福斯蒂努斯主教，以及我们的同司铎斐理伯和阿塞鲁斯在备忘录中带来的那些条文，若在其中有之，则可由我们予以确认；若没有，我们将召开主教会议进一步审议。公证员达尼尔向非洲会议宣读了尼西亚会议的信德宣言及其法令。\par

同样，尼西亚会议的二十条法令也被宣读，正如前面所写的那样。随后，在非洲主教会议中所颁布的事宜被载入当前的会议记录中。\par

% structure: p0049 source=h3 target=subsection reason=relative-heading-rank; base=h2

\subsection{第一条}

\Emph{应严格遵守尼西亚会议的法令}\par

主教奥勒留说：这就是尼西亚会议的法令，我们的教父们当时将其带回；保存此形式，让我们所采纳并确认的以下事宜得以坚定持守。\par

% structure: p0052 source=h3 target=subsection reason=relative-heading-rank; base=h2

\subsection{第二条}

\Emph{论宣讲天主圣三}\par

全体会议说：赖天主的恩宠，在这荣耀的集会中，应首先以一致宣认来宣明我们众所传递的教会信德；然后，各人的教会秩序应通过大家的同意得以建立并加强。为了坚固我们那些新近被祝圣的弟兄和同僚主教们的心灵，应提出我们确实从教父们所领受的，即天主圣三的合一——圣父、圣子及圣神，此合一没有差异，正如我们所说，我们将这样训导天主的子民。此外，所有新近被擢升的主教们说：我们公开宣认，我们这样持守，我们这样教导，遵循福音的信德和你们的教导。\par

% structure: p0055 source=h3 target=subsection reason=relative-heading-rank; base=h2

\subsection{第三条}

\Emph{论节欲}\par

主教奥勒留说：在上次主教会议上讨论节欲与贞洁之事时，那三个等级——即主教、司铎和执事——他们因祝圣而以某种纽带与贞洁相连，看来似乎相宜的是，神圣的统治者、天主的司祭以及肋未人，或那些服务于神圣圣事的人，应当完全节欲，从而能以纯洁的心祈求他们向主所求的；如此，宗徒们所教导、古时所持守的，我们也能持守。\par

% structure: p0058 source=h3 target=subsection reason=relative-heading-rank; base=h2

\subsection{第四条}

\Emph{论不同等级应远离其妻子}\par

福斯蒂努斯，皮切努姆省波滕提纳教会主教、罗马教会代表说：看来相宜的是，主教、司铎和执事，或任何施行圣事的人，应持守端庄，远离其妻子。\par

全体主教说：所有服务祭台的人都应当远离所有妇女，以保持贞洁。\par

% structure: p0062 source=h3 target=subsection reason=relative-heading-rank; base=h2

\subsection{第五条}

\Emph{论贪婪}\par

主教奥勒留说：贪婪的贪欲（无人怀疑它是万恶之母）当从此被禁止，以免任何人侵占他人的地界，或为获利而越出教父所定的界限；任何圣职人员绝不可收取任何形式的利息。那些新颁布的晦涩且通常模棱两可的\OriginalTerm{建议}{suggestiones}，经我们审查后将\OriginalTerm{确定其价值}{formam accipiunt}；但对于神圣圣经已明确判定的事，无需再下判决，只需执行已有规定。在平信徒中可指责的事，在圣职人员中更应受到严厉谴责。全体主教会议说：无人能够不受惩罚地违背先知和福音中所说的话。\par

% structure: p0065 source=h3 target=subsection reason=relative-heading-rank; base=h2

\subsection{第六条}

\Emph{圣化圣油不应由司铎制作}\par

\Person{福尔图纳图斯}主教说：在先前会议中，我们记得曾规定，圣化圣油、忏悔者的和好，以及童贞女的奉献，不应由司铎施行；但若有人被发现如此行了，我们应就他定何规定？\par

\Person{奥勒留}主教说：阁下已听了我们的兄弟和同僚\Person{福尔图纳图斯}主教的建议；你们将作何答复？\par

所有主教答复：圣化圣油的制作与童贞女的奉献，均不应由司铎施行，亦不许司铎在\OriginalTerm{公开弥撒}{publica missa}中与任何人行和好礼，这是众人之共识。\par

% structure: p0070 source=h3 target=subsection reason=relative-heading-rank; base=h2

\subsection{第七条法令}

\Emph{关于那些在死亡危险中获和好者}\par

\Person{奥勒留}主教说：若有人在主教缺席时陷入死亡危险，并寻求与天主的祭台和好，司铎应咨询主教，然后遵其命令为病者行和好礼，此事我们应以健全的谕令加以强化。所有主教说：阁下你所教导我们必要的一切，我等皆赞同。\par

% structure: p0073 source=h3 target=subsection reason=relative-heading-rank; base=h2

\subsection{第八条法令}

\Emph{关于那些控告长上的人；且不得允许罪犯对主教提出指控}\par

\Person{努米底乌斯}，\Place{马克思拉}的主教，说：此外，有许多生活不良的人，认为他们的长上或主教应成为他们控告的靶子；这样的人是否容易准予？\Person{奥勒留}主教说：你们的爱德是否同意，被各样恶行缠累者不得有控告这些人的声诉？\par

所有主教说：若他是有罪之人，他的控告不应被接纳。\par

% structure: p0077 source=h3 target=subsection reason=relative-heading-rank; base=h2

\subsection{第九条法令}

\Emph{关于那些因其行为而被公正逐出教会会众的人}\par

\Person{奥斯定}主教，努米底亚省的代表，说：请明令，若有人因其罪行正当地被逐出教会，后却被某主教或司铎接纳领受共融，则此人应与那些逃避自己主教审判的人同罪。所有主教说：这是众人之共识。\par

% structure: p0080 source=h3 target=subsection reason=relative-heading-rank; base=h2

\subsection{第十条法令}

\Emph{关于被自己主教纠正的司铎}\par

\Person{阿利皮乌斯}主教，努米底亚省的代表，说：也不应忽略此点；若偶有司铎被其主教纠正后，因自负或骄傲而发怒，认为应单独向天主献祭（即脱离主教的权柄），或认为有权另立祭台，违反教会信德与纪律，此人不得免罚。\Person{瓦伦廷}，努米底亚省的首席主教座，说：我们的兄弟\Person{阿利皮乌斯}所提之建议，必然符合教会的纪律与信德；因此，请按诸位的爱德所喜悦的予以颁布。\par

% structure: p0083 source=h3 target=subsection reason=relative-heading-rank; base=h2

\subsection{第十一条法令}

\Emph{若任何司铎因反对其主教而傲慢行事，制造分裂，应受\OriginalTerm{绝罚}{anathema}}\par

所有主教说：若任何司铎被其长上纠正，他应请求邻近主教，使他的案件由他们审理，并经由他们与自己的主教和好；但若他未如此行，反而充满骄傲（愿天主禁止！），认为应自行脱离与主教的共融，并单独向天主献祭，且与某些同谋者制造分裂，则他应受\OriginalTerm{绝罚}{anathema}，并丧失其职位；若他所提出对其主教的控告经查并非属实，则应进行调查。\par

% structure: p0086 source=h3 target=subsection reason=relative-heading-rank; base=h2

\subsection{第十二条法令}

\Emph{若任何主教在会议期间外受到控告，其案件应由十二位主教审理}\par

\Person{斐理克斯}主教说：我提议，依据古代会议的规定，若任何主教（愿天主禁止！）受到某项控告，且过于紧急无法等待召集多数主教，为使他不致长期受指控，其案件应由十二位主教审理；司铎应由六位主教连同其本主教审理，执事应由三位审理。\par

% structure: p0089 source=h3 target=subsection reason=relative-heading-rank; base=h2

\subsection{第十三条法令}

\Emph{主教不应由多人之外祝圣，但若有需要，可由三位主教祝圣}\par

\Person{奥勒留}主教说：诸位圣德对此有何意见？所有主教答道：古人所定之规定，我们必须遵守，即没有教省首席之同意，即使许多主教聚集，也不得轻率祝圣主教。但若有需要，遵其命令，可由三位主教在任何所在祝圣他；若有人违反其誓言与签名而进入某地，他将因此丧失其尊位。\par

% structure: p0092 source=h3 target=subsection reason=relative-heading-rank; base=h2

\subsection{第十四条法令}

\Emph{\Place{的黎波里}的一位主教应作为代表前来，且在那里司铎可由五位主教审理}\par

此外，也以为合宜：因该省贫困，可由\Place{的黎波里}的一位主教作为代表前来，且在那里，如前所述，司铎应由五位主教审理，执事由三位，其本主教主持。\par

% structure: p0095 source=h3 target=subsection reason=relative-heading-rank; base=h2

\subsection{第十五条法令}

\Emph{论服务于教会的各种品位：若任何人陷入刑事事务并拒绝由教会法庭审判，他应因此处于危险；且主教（\OriginalTerm{司铎的}{sacerdotum}）的儿子们不得参加世俗表演。}\par

此外，以为合宜：若任何主教、司铎或执事被提起刑事控告，或有民事案件，拒绝由教会法庭审判，而愿由世俗法庭裁决，即使他胜诉，仍应丧失其职位。\par

此乃刑事案件的法规；但在民事案件中，他若愿保留其职位，则应丧失其所起诉争取之物的权益。\par

此外，也认为为好，倘若有人从某些教会法官向上级管辖权的其他教会法官提起上诉，这决不应损害被上诉之人的声誉，除非能证明他们出于仇恨或其他某种精神偏见而作出判决，或因某种方式被腐化。但若经双方同意选出了法官，即使人数少于规定，也不得上诉。\par

并且〔认为为好〕，主教之子不应参与或观看世俗的表演。因为这事向来禁止所有基督徒，所以让他们远离这些，以免前往有咒骂和亵渎之处。\par

% structure: p0101 source=h3 target=subsection reason=relative-heading-rank; base=h2

\subsection{第十六条。}

\Emph{主教、司铎或执事不得担任\OriginalTerm{管理人}{conductor}；读经员应娶妻；神职人员应戒除高利贷；以及他们在什么年龄应被祝圣，或童贞女在什么年龄应被祝圣。}\par

同样，认为为好，主教、司铎和执事不得为\OriginalTerm{管理人}{conductor}或\OriginalTerm{代理人}{procurator}；也不得以任何卑贱和恶劣的职业谋生，因为他们应记住经上记载：\Quote{没有人服事天主而让自己纠缠于世俗事务。}\par

再者，认为为好，读经员到了青春期时，应被迫要么娶妻，要么宣示守贞。\par

同样，认为为好，倘若神职人员借钱，应收回本金；但若借出的是\OriginalTerm{种类}{speciem}，则应收回他所给出的同一种类。\par

并且，未满二十五岁的执事不得被祝圣，童贞女也不得被祝圣。\par

并且，读经员不应向人群致意。\par

% structure: p0108 source=h3 target=subsection reason=relative-heading-rank; base=h2

\subsection{第十七条。}

\Emph{任何省份因其距离较远，可拥有自己的首席主教。}\par

认为为好，\Place{毛里塔尼亚·西提芬西斯}，如其请求，应有自己的首席主教，经\Place{努米底亚}首席主教同意，该省已从其教务会议中分离。并且经\Place{非洲各省}所有首席主教及所有主教同意，鉴于彼此相距甚远，给予许可。\par

% structure: p0111 source=h3 target=subsection reason=relative-heading-rank; base=h2

\subsection{第十八条。（希腊文第十八条。\Emph{拉丁文标题为希腊文教规。}）}

\Emph{若有任何神职人员被祝圣，应劝告他遵守各项规定。}\par

并且，圣体圣事和圣洗圣事均不应给予死者的遗体。\par

并且，每年在各省，都主教应召开教务会议。\par

% structure: p0115 source=h3 target=subsection reason=relative-heading-rank; base=h2

\subsection{（希腊文第十九条。）}

认为为好，在主教或神职人员被祝圣之前，应将教规的条文知会他们，以免他们日后因无知而违法，后悔莫及。\par

% structure: p0117 source=h3 target=subsection reason=relative-heading-rank; base=h2

\subsection{（希腊文第二十条。）}

同样认为为好，圣体圣事不应给予死者的遗体。因为经上记载：\Quote{拿去，吃。}但死者的遗体既不能\Quote{拿去}也不能\Quote{吃。}也不可因司铎的无知而为死者施洗。\par

% structure: p0119 source=h3 target=subsection reason=relative-heading-rank; base=h2

\subsection{（希腊文第二十一条。）}

因此，在此神圣教务会议上，应依照尼西亚法令予以确认，鉴于教会事务常常因拖延而损害人民利益，每年应举行教务会议，各省首席主教应派遣主教作为代表，从其省内教务会议中选出两名或更多，以便教务会议召开时拥有充分的行动权。\par

% structure: p0121 source=h3 target=subsection reason=relative-heading-rank; base=h2

\subsection{第十九条。（希腊文第二十二条。）}

\Emph{若有主教被控告，案件应提交其本省首席主教。}\par

\Person{奥勒留}主教说：凡被控告的主教，原告应将案件提交被控告者所属省的首席主教，被告不应因被指控的罪行而被停职共融，除非他未能在指定日期出庭，在由信函正式传唤后，在选定的法官面前辩护其案件；即自他确认收到信函之日起一个月内。但若能证明有任何真正的必要情况，显然使其无法出庭，他应有机会再获整整一个月来辩护其案件；但第二个月过后，在宣判无罪之前，他不得参与共融。\par

但若他不愿参加年度总会议，以使他的案件在那里终结，他应被视为在自己不共融之时宣判了对自己定罪的判决，并且他不得在自己的教会或教区共融。\par

但若原告未曾错过任何辩护日期，则不应被排除共融；但若他错过了其中一些日期，自行退避，那么主教应恢复共融，而原告应被移除共融；然而，仍保留他继续诉讼的可能性，若他能证明自己的缺席是因无力而非不愿意。\par

并且为此规定，若在针对主教的案件进入辩论时，原告本身被证明是罪犯，则不准其作证，除非他声称事由是个人性的而非教会性的。\par

% structure: p0127 source=h3 target=subsection reason=relative-heading-rank; base=h2

\subsection{第二十条。（希腊文第二十三条。）}

\Emph{论被控告的司铎或神职人员。}\par

但若司铎或执事被控告，应从邻近地方与教区主教一同召集法定数量的同僚，由被告自行从同一范围指定；即针对司铎的案件中，连同被告本人共六人；针对执事的案件中，三人。他们应讨论案件，并且关于日期、延期和移除共融，以及关于原告和被告之间对人员的讨论，应遵守相同的形式。\par

至于其余神职人员的案件，当地主教应单独审理并裁决。\par

% structure: p0131 source=h3 target=subsection reason=relative-heading-rank; base=h2

\subsection{第二十一条。（希腊文第二十四条。）}

\Emph{神职人员之子不得与异端者联姻。}\par

同样，认为为好，神职人员之子不应与外邦人和异端者缔结婚姻。\par

% structure: p0134 source=h3 target=subsection reason=relative-heading-rank; base=h2

\subsection{第二十二条。（希腊文第二十五条。）}

\Emph{主教或其他神职人员不得将任何东西给予非天主教徒。}\par

并且，对非天主教徒，即使是血亲，主教或神职人员也不得将其财产以捐赠方式给予任何东西。\par

% structure: p0137 source=h3 target=subsection reason=relative-heading-rank; base=h2

\subsection{第二十三号法令（希腊编号二十六）}

\Emph{主教不得渡海}\par

再者，主教若无本省首席主教同意，不得渡海；以便从该首席主教获得格式信函或推荐信。\par

% structure: p0140 source=h3 target=subsection reason=relative-heading-rank; base=h2

\subsection{第二十四号法令（希腊编号二十七）}

\Emph{除正典圣经外，教会中不得宣读任何其他读物}\par

再者，除正典圣经外，不得以神圣圣经之名在教会中宣读任何其他读物。\par

正典圣经如下：\par

\begin{itemize}
\item \WorkTitle{创世纪}。
\end{itemize}

\begin{itemize}
\item \WorkTitle{出谷纪}。
\end{itemize}

\begin{itemize}
\item \WorkTitle{肋未纪}。
\end{itemize}

\begin{itemize}
\item \WorkTitle{户籍纪}。
\end{itemize}

\begin{itemize}
\item \WorkTitle{申命纪}。
\end{itemize}

\begin{itemize}
\item \WorkTitle{若苏厄书}。
\end{itemize}

\begin{itemize}
\item \WorkTitle{民长纪}。
\end{itemize}

\begin{itemize}
\item \WorkTitle{卢德传}。
\end{itemize}

\begin{itemize}
\item \WorkTitle{列王纪，四卷}。
\end{itemize}

\begin{itemize}
\item \WorkTitle{编年纪，两卷}。
\end{itemize}

\begin{itemize}
\item \WorkTitle{约伯传}。
\end{itemize}

\begin{itemize}
\item \WorkTitle{圣咏集}。
\end{itemize}

\begin{itemize}
\item \WorkTitle{撒罗满五书}。
\end{itemize}

\begin{itemize}
\item \WorkTitle{十二小先知书}。
\end{itemize}

\begin{itemize}
\item \WorkTitle{依撒意亚}。
\end{itemize}

\begin{itemize}
\item \WorkTitle{耶肋米亚}。
\end{itemize}

\begin{itemize}
\item \WorkTitle{厄则克耳}。
\end{itemize}

\begin{itemize}
\item \WorkTitle{达尼尔}。
\end{itemize}

\begin{itemize}
\item \WorkTitle{多俾亚传}。
\end{itemize}

\begin{itemize}
\item \WorkTitle{友弟德传}。
\end{itemize}

\begin{itemize}
\item \WorkTitle{艾斯德尔传}。
\end{itemize}

\begin{itemize}
\item \WorkTitle{厄斯德拉，两卷}。
\end{itemize}

\begin{itemize}
\item \WorkTitle{玛加伯，两卷}。
\end{itemize}

\begin{itemize}
\item \WorkTitle{新约}。
\end{itemize}

\begin{itemize}
\item \WorkTitle{福音书，四卷}。
\end{itemize}

\begin{itemize}
\item \WorkTitle{宗徒大事录，一卷}。
\end{itemize}

\begin{itemize}
\item \WorkTitle{保禄书信，十四封}。
\end{itemize}

\begin{itemize}
\item \WorkTitle{伯多禄宗徒书信，两封}。
\end{itemize}

\begin{itemize}
\item \WorkTitle{若望宗徒书信，三封}。
\end{itemize}

\begin{itemize}
\item \WorkTitle{雅各伯宗徒书信，一封}。
\end{itemize}

\begin{itemize}
\item \WorkTitle{犹达宗徒书信，一封}。
\end{itemize}

\begin{itemize}
\item \WorkTitle{若望默示录，一卷}。
\end{itemize}

应将此送交我们的弟兄和同为主教的\Person{博尼法斯}，以及那些地区的其他主教，使他们确认此法令，因为这些是从我们祖先所领受、应在教会中宣读的经书。\par

% structure: p0177 source=h3 target=subsection reason=relative-heading-rank; base=h2

\subsection{第二十五号法令（希腊编号二十八）}

\Emph{论主教及侍奉至圣奥迹的较低品级。决定这些人员应戒绝其妻子}\par

\Person{奥勒留}主教说：我们还要补充，最亲爱的弟兄们，既然我们听说某些神职人员，甚至包括读经员，对其妻子行为不节制，故此决定：此前各会议所颁布的法规应予确认，即侍奉神圣奥迹的副执事、执事、司铎及主教，应依照先前的规定，与妻子分房，如同无妻一般；若不如此行，应撤其职务。但其余神职人员，除非年事已长，不必强制。全体会议说道：您圣洁所言公正、神圣且悦乐天主，我们予以确认。\par

% structure: p0180 source=h3 target=subsection reason=relative-heading-rank; base=h2

\subsection{第二十六号法令（希腊编号二十九）}

\Emph{无人应取教会的财产。}\par

同样，决定：无人应出售属于教会之物；若无收益，且另有重大急需催迫，应将此事提交本省都主教，由其与指定数目之主教商议是否可行；若某教会遇如此紧急需要，以致无法事先商议，至少应召集邻近主教作为见证，并注意将其教会之一切需要提交议会；若他不如此行，则在天主面前被视为有责，在议会眼中被视为售卖者，并因此失去其荣誉。\par

% structure: p0183 source=h3 target=subsection reason=relative-heading-rank; base=h2

\subsection{第二十七号法令（希腊编号三十）}

\Emph{司铎和执事被判定犯有较严重罪行时，不得如平信徒般接受按手礼。}\par

亦经确认：若司铎或执事因犯有重大罪行而须被撤职，对他们不得如同对忏悔者或信友平信徒那样行按手礼，也不得允许他们重新受洗，然后晋升至神职品级。\par

% structure: p0186 source=h3 target=subsection reason=relative-heading-rank; base=h2

\subsection{第二十八号法令（希腊编号三十一）}

\Emph{司铎、执事或神职人员，若决意将其案件上诉至海外，绝不得被接纳共融。}\par

亦决定：司铎、执事及其他较低级神职人员，如对其主教的判决不满，可在其主教同意下由邻近主教审理，被召来的主教应在他们之间裁决；但若他们认为有理由对此上诉，不得诉诸海外的审判，而应诉诸本省的首席主教，或至普世大公会议，如同关于主教的法令也如此规定。但任何决意将上诉带到海外者，在非洲境内不得被任何人接纳共融。\par

% structure: p0189 source=h3 target=subsection reason=relative-heading-rank; base=h2

\subsection{第二十九号法令（希腊编号三十二）}

\Emph{凡被绝罚者，若在其案件听证前擅自领受共融，则自招惩罚。}\par

同样，全体会议悦纳：凡因任何疏忽被绝罚者，无论是主教或其他神职人员，若在判决未撤销、案件未听证时擅自领受共融，则此举应被视为自我定罪。\par

% structure: p0192 source=h3 target=subsection reason=relative-heading-rank; base=h2

\subsection{第三十号法令（希腊编号三十三）}

\Emph{关于被告或原告。}\par

同样，决定：被告或原告，若（与被控者同住一处）担心恶人暴民加害，可为自己另选近处，以便案件得以审断，且易于传唤证人。\par

% structure: p0195 source=h3 target=subsection reason=relative-heading-rank; base=h2

\subsection{第三十一号法令（希腊编号三十四）}

\Emph{若某些神职人员被其主教擢升后傲慢自负，则不得留于其不愿离去之处。}\par

亦决定：凡神职人员或执事，在主教希望将其提拔至教区中更高职位时，若不愿在教会急需中协助主教，则不得再执行其不愿被调离之品级的职务。\par

% structure: p0198 source=h3 target=subsection reason=relative-heading-rank; base=h2

\subsection{第三十二号法令（希腊编号三十五）}

\Emph{若任何贫穷神职人员，无论其品级如何，获得任何财产，该财产应归主教管辖。}\par

众主教、司铎、执事及其他圣职人员，凡在受圣职时本无财产，而在任主教期间或成为圣职人员之后，以其本人名义购置土地或其他财产者，应被视为侵占主之财物之罪，除非在受劝告后，将此等财产归于教会。但若有人因他人慷慨赠予或亲属继承而获得财产，则任凭其处置；然而，若他们反悔而索回，则将被视为退后者，不配享有教会之尊荣。\par

% structure: p0201 source=h3 target=subsection reason=relative-heading-rank; base=h2

\subsection{第33条（希腊文第36条）}

\Emph{司铎不得出售其就职之教会的财产，任何主教不得正当地使用任何产权属于教会母中心的物品} (\OriginalTerm{母中心}{\GreekText{μάτρικος}})。\par

主教非因必要，不得侵占\Emph{母的} (\OriginalTerm{母堂}{matricis}) 教会的财产 [司铎也不得侵占其本堂 (\OriginalTerm{本堂}{tituli}) 的财产]。\par

% structure: p0204 source=h3 target=subsection reason=relative-heading-rank; base=h2

\subsection{第34条（希腊文第37条）}

\Emph{希波主教会议所制定的诸项，不得修改任何一条}\par

\Person{埃庇戈尼乌斯}主教说：在此\OriginalTerm{纲要}{Breviarium}中，我们认为不应作任何修改，也不应添加任何内容，除非在主教会议上宣布圣逾越节之日期。\par

% structure: p0207 source=h3 target=subsection reason=relative-heading-rank; base=h2

\subsection{第35条（希腊文第38条）}

\Emph{主教或圣职人员不应轻易释放其子女}\par

主教或圣职人员不应轻易让子女脱离其权力，除非对他们的品行和年龄有把握，以致他们自身的罪过只归咎于自己。\par

% structure: p0210 source=h3 target=subsection reason=relative-heading-rank; base=h2

\subsection{第36条（希腊文第39条）}

\Emph{除非使全家成为基督徒，否则不得祝圣主教或圣职人员}\par

任何人不得被祝圣为主教、司铎或执事，除非其家中所有成员都已成为公教徒基督徒。\par

% structure: p0213 source=h3 target=subsection reason=relative-heading-rank; base=h2

\subsection{第37条（希腊文第40条）}

\Emph{在神圣奥秘中，除搀水之饼与酒外，不得奉献任何物品}\par

在主的身体与血的圣事中，只应奉献主所亲自规定的，即搀水之饼与酒。至于初熟之物，无论蜂蜜或牛奶，可在最为隆重的那一日，按照惯例，在婴儿的奥秘中奉献。虽然它们被奉献在祭台上，但仍应有其独特的祝福，以区别于主之身体与血的圣事；初熟之物不得奉献葡萄和谷穗以外的物品。\par

% structure: p0216 source=h3 target=subsection reason=relative-heading-rank; base=h2

\subsection{第38条（希腊文第41条）}

\Emph{圣职人员或守贞者不得访问贞女或寡妇}\par

圣职人员或宣示守贞者，未经主教或司铎的许可或同意，不得进入寡妇或贞女的住所；即使进入，也不得独自前往，而应有其他圣职人员或主教、司铎为此指定的人陪同；主教和司铎也不得单独访问此类妇女，除非有其他圣职人员或庄重的基督徒男子在场。\par

% structure: p0219 source=h3 target=subsection reason=relative-heading-rank; base=h2

\subsection{第39条（希腊文第42条）}

\Emph{主教不应被称为司祭之首}\par

首席主教不应被称为司祭之王或\OriginalTerm{大司祭}{Summus Sacerdos}，或任何其他此类称号，而仅应被称为首席主教。\par

% structure: p0222 source=h3 target=subsection reason=relative-heading-rank; base=h2

\subsection{第40条（希腊文第43条）}

\Emph{关于圣职人员不得出入酒馆，除旅行外}\par

圣职人员不得进入酒馆进食或饮酒，除非因旅行需要而不得已为之。\par

% structure: p0225 source=h3 target=subsection reason=relative-heading-rank; base=h2

\subsection{第41条（希腊文第44条）}

\Emph{祭献应由守斋者向天主呈奉}\par

祭台上的圣事不应由非守斋者举行，除非在纪念主的晚餐的周年庆典上；因为若要在午后为某些死者（无论是主教或其他）举行纪念，则只可进行祈祷，而主礼者若已用过早餐。\par

% structure: p0228 source=h3 target=subsection reason=relative-heading-rank; base=h2

\subsection{第42条（希腊文第15条）}

\Emph{在任何情况下不得在教堂内举行宴饮}\par

任何主教或圣职人员不得在教堂内举行宴饮，除非因行路之需而被迫招待旅客。也应尽可能禁止民众参与此类宴饮。\par

% structure: p0231 source=h3 target=subsection reason=relative-heading-rank; base=h2

\subsection{第43条（希腊文第46条）}

\Emph{关于补赎者}\par

补赎者的补赎期限应由主教视其罪过之不同而定；司铎未经请示主教，不得与补赎者和好，除非因主教缺席而不得已为之。但若补赎者的过犯是公开且众所周知的，以致令整个教会蒙受恶表，则应在祭台前 (拉丁文\Quote{半圆形后殿前}) 接受覆手礼。\par

% structure: p0234 source=h3 target=subsection reason=relative-heading-rank; base=h2

\subsection{第44条（希腊文第47条）}

\Emph{关于贞女}\par

圣洁的贞女在离开平时监护她们的父母后，应由主教或（主教不在时）司铎托付给年长的妇女，与她们同住，以便照顾她们，免得她们四处游荡，损害教会的声誉。\par

% structure: p0237 source=h3 target=subsection reason=relative-heading-rank; base=h2

\subsection{第45条（希腊文第48条）}

\Emph{关于患病而不能自行回答者}\par

患病而不能自行回答者，若其[仆人]以自己的危险为[病人]的善意作证，则可为之施洗。\par

（希腊文第49条）\par

\Emph{关于正在补赎并归向主的优伶}\par

对于优伶、演员及其他此类人士，以及背教者，当他们悔改并归向天主时，不应拒绝恩宠或和好。\par

% structure: p0243 source=h3 target=subsection reason=relative-heading-rank; base=h2

\subsection{第46条（希腊文第50条）}

\Emph{关于殉道者的苦难}\par

当庆祝殉道者的周年纪念日时，可诵读殉道者的受难录。\par

% structure: p0246 source=h3 target=subsection reason=relative-heading-rank; base=h2

\subsection{第四十七法令。（希腊文第五十一则）}

\Emph{关于[多纳徒派和]受多纳徒派洗礼的孩童。}\par

关于多纳徒派，我们认为应与我们的弟兄和同司铎西利奇乌斯及辛普利西安商议，仅针对那些独自受多纳徒派洗礼的婴孩；以免当他们以得救的决心归向天主的教会时，他们父母之错误可能阻碍他们被提升至圣祭台的职务，因为他们并未出于自己的意愿行事。\par

但当这些事开始进行时，西提芬毛里塔尼亚省的主教奥诺拉图斯和乌尔巴努斯说：前些日子我们被派往贵圣座时，我们搁置了为此事所写的内容，以便等候来自努米底亚的弟兄特使到达。但因已过了不少日子，我们等候他们却仍未到达，故此不宜再延迟我们接受自同为主教的弟兄们的命令；所以，弟兄们，请以欣然之心接纳我们的陈述。我们听闻关于尼西亚信经的信德：确实，早餐后应禁止祭献，好让守斋者能恰当地奉献，这一规定在当时和现在都已得到确认。\par

% structure: p0250 source=h3 target=subsection reason=relative-heading-rank; base=h2

\subsection{第四十八法令。（希腊文第五十二则）}

\Emph{论重洗、重覆祝圣及迁移主教}\par

但我们建议，我们应颁令卡普亚全体会议以其智慧所陈述的：不得举行重洗、不得重覆祝圣，亦不得迁移主教。因为维拉雷吉斯主教克雷斯肯提乌斯离开自己的羊群，侵占了图比尼亚教会，且至今受劝告要他按照法令离开他所侵占的教区，他却藐视这劝告。我们听闻对他所宣判的判决已获确认；但根据我们的法令，我们请求你们惠允，在必要之下，我们可向辖省长官申诉，依照最荣耀的君主的律法，以便凡不愿顺从贵圣座的温和劝告并改正其不法行为者，当立即被司法权威逐出。奥雷利乌斯主教说：依照既定规则的遵守，如果他被你们，亲爱的弟兄们，要求离开而拒绝，则不得被视为会议的一员；因为他因自身的藐视和悖逆，已落入世俗法官的权力之下。奥诺拉图斯和乌尔巴努斯主教说：这岂不令我们众人悦纳？众主教答说：这是正义的，我们悦纳。\par

% structure: p0253 source=h3 target=subsection reason=relative-heading-rank; base=h2

\subsection{第四十九法令。（希腊文第五十三则）}

\Emph{视圣一位主教当有多少主教在场}\par

\Person{奥诺拉图斯}和\Person{乌尔巴努斯}主教说：我们颁令此命，因为最近我们两位努米底亚的主教弟兄擅自祝圣了一位主教，故主教祝圣仅当有十二位主教共聚时举行。奥雷利乌斯主教说：应当保存古老规则，即指定不少于三位主教足以祝圣主教。何况，在的黎波里和阿尔祖格，蛮族如此邻近，因声称在的黎波里仅有五位主教，其中两位或因某种需要不得空闲；然而所有主教齐聚一处实属困难；难道这情况应妨碍完成对教会有益的事么？因为在此教会中，贵圣座已惠允召集我们，我们常有祝圣，几乎每个主日；我岂能经常召集十二位、十位或约略此数的主教？但对我而言，联合两三位邻近主教与我卑微之身，是容易之事。因此，你们的爱德当同意我，此规条不能遵守。\par

% structure: p0256 source=h3 target=subsection reason=relative-heading-rank; base=h2

\subsection{第五十法令。（希腊文第五十四则）}

\Emph{若有人反对被祝圣者，当加添多少主教于祝圣者数目中}\par

但应规定：当我们聚集拣选一位主教时，若发生任何反对，因我们已处理过此类事件，那三位主教不得擅自洗清将被祝圣者，而应请求在前述数目之外再加一两位主教，且先应在同一地方（\OriginalTerm{会众}{plebe}）中审查反对者的人。最后再审查反对理由；他仅在公众面前被洗清后方可予以祝圣。若此提议合乎贵圣座之意，请以贵尊驾的回答予以确认。众主教说：我们甚为满意。\par

% structure: p0259 source=h3 target=subsection reason=relative-heading-rank; base=h2

\subsection{第五十一法令。（希腊文第五十五则）}

\Emph{复活节日期应由迦太基教会公布}\par

\Person{奥诺拉图斯}和\Person{乌尔巴努斯}主教说：既然我们的备忘录所处理的一切都已明了，我们也补充关于复活节日子所吩咐的：我们始终由迦太基教会获知日期，一如惯常，且不可太迟。奥雷利乌斯主教说：若贵圣座认为合宜，因我们记得此前曾承诺每年聚集讨论，当我们集会时，当由出席大公会议的特使公布圣复活节的日子。\Person{奥诺拉图斯}和\Person{乌尔巴努斯}主教说：现在我们请求本次集会，惠允以信函通知我们省份那日子。奥雷利乌斯主教说：理当如此。\par

% structure: p0262 source=h3 target=subsection reason=relative-heading-rank; base=h2

\subsection{第五十二法令。（希腊文第五十六则）}

\Emph{论巡视各省}\par

\Person{奥诺拉图斯}和\Person{乌尔巴努斯}主教说：曾口头命令我们，因希波会议曾规定，各教省应在会议期间被巡视，也请俯允今年或明年，按照您所制定的顺序，巡视毛里塔尼亚省。\par

奥勒留主教说：关于毛里塔尼亚行省，因它位于阿非利加的边界，我们未作任何规定，因他们与蛮族相邻；但愿天主恩赐（然而我并非轻率承诺必能做到），使我们能前往你们的行省。因为你们应思考，弟兄们，的黎波里和阿尔祖格斯地区的我们这些弟兄，若有适当机会，也会提出同样的要求。\par

% structure: p0266 source=h3 target=subsection reason=relative-heading-rank; base=h2

\subsection{第五十三号教规（希腊文本第五十七号）}

教区未经其所属主教同意，不得另受主教\par

\Person{埃皮戈尼乌斯}主教说：许多次会议中，司铎会议已规定，凡属其他教区并由其主教管辖、且从未有过自己主教的团体，未经管辖他们的主教同意，不得为自己接受统治者，即主教。但因某些获得某种统治权的人，厌恶弟兄的共融，或至少变得堕落，以真正的暴政为自己要求统治权，大多是傲慢愚蠢的长老，他们抬头反对自己的主教，或通过宴请或恶意的说服赢得民众，好藉不法恩宠将自己立为他们的统治者。我们实在坚持您心灵的荣光，最敬爱的弟兄奥勒留，因您常反对这些事，不理会这些请愿者；但鉴于他们邪恶的念头和卑劣的阴谋，我说，这样的团体，既然始终隶属某一教区，就不应接受一位统治者，也永远不应拥有自己的主教。因此，若我所提议的为整个至圣会议所认可，就请予以确认。\par

奥勒留主教说：我不反对我们弟兄和同僚主教的提议；但我承认，对于真正同心合意的人，不仅关乎迦太基教会，也关乎每一次司铎集会，我一向如此行，将来也如此行。因为有许多人，正如所说，与他们所欺骗的民众共谋，搔他们的耳朵，谄媚地引诱他们，这些人品行恶劣，或至少自高自大，与本次会议分离，以为自己能照管自己的民众，却从不来参加我们的会议，恐怕他们的恶行被讨论。我说，若认为合适，这些人不仅不应保留他们的教区，还应竭力通过公共权力将他们从那个不当地偏袒他们的教会中驱逐出去，甚至应将其从首席教座上移除。因为那与所有弟兄及整个会议联合的人，应充分享有管理自己教会及教区的权利。但那些以为民众就足够、而蔑视弟兄之爱的人，不仅将失去他们的教区，而且（如我所说）应被公共权力剥夺他们自己的牧职，作为反叛者。\Person{霍诺拉图斯}和\Person{乌尔班}主教说：您圣座的高瞻远瞩获得我们全体心灵的认同，我认为，通过大家的回答，您所屈尊提议的将得到确认。所有主教说：赞成，赞成。\par

% structure: p0270 source=h3 target=subsection reason=relative-heading-rank; base=h2

\subsection{第五十四号教规（希腊文本第五十八号）}

任何情况下，不得接纳外方圣职人员\par

\Person{埃皮戈尼乌斯}主教说：许多次会议中已有此规定，且刚才又蒙诸位至福弟兄的明智所确认：任何主教不得在未经该圣职人员所属主教的同意下，将其外方圣职人员接纳进入自己的教区。但我说，有个名叫犹利安的人，对天主通过小可赐予他的恩赐忘恩负义，竟如此鲁莽大胆，以致一个曾由我施洗的人——当他还是一名极贫困的少年时，由他本人托付给我，且多年由我喂养抚养——此人确如我所说，在我自己的教堂中由我卑微的双手受洗；他在马帕连教区开始担任读经员，在那里读了近两年，极其藐视小可。前述犹利安将此声称是他自己城市\Place{瓦扎里塔}的公民之人，未与我商议就祝圣他为执事。至福的弟兄们，若此事可容许，请告诉我们；若不可容许，就应约束如此厚颜无耻之人，使他绝不再混入某人的共融之中。\par

\Person{努米底乌斯}主教说：若如所看来，犹利安未经向您尊驾请求同意，甚至未曾咨询，就做了此事，我们全体判定此事做得不义且不值。因此，除非犹利安纠正其错误，并以适当的赔补将那位圣职人员归还您的子民，鉴于他所做违反会议规定，就因其顽抗应被定罪，与我们分离。\Person{埃皮戈尼乌斯}主教说：我们年长的父亲，在晋升上最古老的、可赞美的弟兄和同僚\Person{维克托}希望此请求能普及众人。\par

% structure: p0274 source=h3 target=subsection reason=relative-heading-rank; base=h2

\subsection{第五十五号教规（希腊文本第五十九号）}

迦太基主教得随己意祝圣圣职人员\par

奥勒留主教说：我的弟兄们，请允许我发言。常有身处需要的圣职人员向我寻求执事（拉丁文为\OriginalTerm{负责人}{præpositis}），或长老，或主教。我牢记所规定的事，遵循如下：我召集被寻求的圣职人员所属的主教，向他说明情况，某教会请求他某位圣职人员。或许他们不反对，但恐怕后来当他们已被我（你们知道，我负责许多教会及受祝圣者）要求时，他们又反对。因此，我该召集一位同僚主教，并有两三位我们中的证人。但若发现他\OriginalTerm{不虔诚的}{\GreekText{ἀκαθοσίωτος}}，你们的爱德认为应如何做？因为你们知道，弟兄们，因天主的俯就，我负责所有教会。\par

\Person{Numidius}主教说：\Quote{此座一直有权依照各教会的心愿，按自己的意愿，并在众人一致同意（\Emph{fuisset conventus}）的人上祝圣主教。}\Person{Epigonius}主教说：\Quote{您善良的天性很少运用您的权力，因为您使用它们远少于您所能使用；因为，我的兄弟，您对所有人都仁慈和蔼；您有权柄，但就仅在\Emph{prima tantummodo conventione}上满足每位主教的个人意愿，这远非您的习惯。但如果认为此座的权利应当得到主张，那么您有责任支持所有教会，因此我们并非赋予您权力，而是确认您已有的权力，即：您有权始终按自己的意愿选择任何人，在那些请求您的人民和教会中立他们为首牧，只要您愿意。}\Person{Posthumianus}主教说：\Quote{一个只有一位司铎的人，要把他那一位带走，这合理吗？}\Person{Aurelius}主教说：\Quote{但可能有一位主教，通过上主的恩惠能使许多人成为司铎，但适合成为主教的人却很难找到。因此，如果某位主教有一位对主教职必要的司铎，而他只有那一人，我的兄弟，如您所说，即便是那一位，他也应当为晋升而奉献出来。}\Person{Posthumianus}主教说：\Quote{如果其他某位主教有充足的圣职人员，那么那个别的教区应该来帮助我吗？}\Person{Aurelius}主教说：\Quote{当然，当您帮助了另一教会时，拥有许多圣职人员的人应当被说服，让出一位给您祝圣。}\par

% structure: p0278 source=h3 target=subsection reason=relative-heading-rank; base=h2

\subsection{第56条（希腊本第六十条）}

\Emph{被祝圣为教区的主教不得为自己选择教区［在希腊各省］。}\par

\Person{Honoratus}和\Person{Urban}两位主教说：\Quote{我们听说已经规定，除非得到建立者的同意，教区不应被认为适合接受主教；但在我们的省份，由于有些人已经在教区被祝圣为主教，而且是在建立他们的主教同意下，他们甚至为自己夺取了教区，这应当由您的爱德审断予以纠正，并禁止将来再发生。}\Person{Epigonius}主教说：\Quote{每位主教应有的权利应当保留，以便从教区整体中被剥离出任何部分而拥有自己的主教时，必须获得适当权力者的同意。因为只要给予同意，被分离出来的教区仅拥有自己的主教就足够了，且不可让他夺取其他教区，因为只有从众多中分出的一块才配得接受主教的荣誉。}\Person{Aurelius}主教说：\Quote{我不怀疑你们各位的爱德都乐于接受，那位为了一教区而被祝圣、且得到拥有母堂的主教的同意者，只应保留他为之祝圣的人民。既然我认为所有事宜都得到了处理，如果一切合你们的心意，请用你们的投票确认全部。}所有主教说：\Quote{我们都很满意，并已通过我们的签署确认。}然后他们签署了名字。\par

\Signature{我，\Person{Aurelius}，迦太基教会的主教，已同意此法令，并签署了所宣读的内容。所有其他主教也都同样签署。}\par

% structure: p0282 source=h3 target=subsection reason=relative-heading-rank; base=h2

\subsection{第57条（希腊本第六十一条）}

\Emph{儿童时期受多纳特派洗礼者，可在天主教会中被祝圣为圣职人员。}\par

由于在先前的大公会议中，如同诸位的同心合意和我一样所记得的，曾规定：那些自幼被多纳徒派受洗的儿童，尚未能认识其错误的危害，后来当他们达到理性年龄，接受了真理的知识，憎恶他们先前的错误，并按照古代惯例，藉着覆手而被接纳进入散佈普世的天主教会时，这样的错误记忆不应妨碍他们接受圣职。因为当他们进入信仰时，他们以为自己的教会就是真教会，并在那里相信了基督，领受了天主圣三的圣事。而这些圣事全都是真实、神圣而属神的，这是最为确定的，并且整个灵魂的希望就寄托在其中，尽管异端者的狂妄僭越，冒用真理之名，竟敢施行它们。然而它们毕竟只是一，正如蒙福的宗徒告诉我们的，说：\Quote{一位天主、一个信德、一个圣洗}，而且不可重复那只应一次施行的事。[因此那些如此受洗的人]在弃绝他们的错误后，可藉着覆手被接纳进入唯一的教会，即被称为柱石和所有基督徒的唯一母亲，在那里所有这些圣事都被领受而获得救恩和永生；甚至同一圣事，对于坚持异端的人，则招致沉重的惩罚。所以，那对于在真理中的人光照他们获得永生的事物，对于在错误中的人反而趋向黑暗和灭亡。因此，对于那些逃离错误、承认他们的母亲——天主教会——的怀抱，以真理的爱相信并领受所有这些神圣奥秘，并且除了圣事之外还有美好生活见证的人，没有人会不同意，这样无疑可以晋升圣职，尤其是在像当前这样的急需情况下。然而，还有这一派别中的其他人，他们已经是神职人员，希望带着他们的教友和他们的荣衔转到我们这里，例如那些为职位的缘故皈依得救的人，并且为了保留这些职衔而寻求救恩[即进入教会]。我认为关于此类人的问题可以留给我们前述弟兄的更审慎的考虑，并且当他们以更明智的顾问考虑提交给他们的事项后，他们可惠允就这一问题告知我们他们认为合适的意见。只是关于那些自幼被异端者受洗的儿童，我们决定：如果他们同意，就赞成我们关于他们祝圣的判决。因此，凡我们以上与神圣主教们所陈述的一切，愿你们的可敬弟兄与我一起判定执行。\par

% structure: p0285 source=h3 target=subsection reason=relative-heading-rank; base=h2

\subsection{第五十八（希腊编号第六十二）条}

关于剩余的偶像或庙宇应由皇帝废除\par

\Emph{因此}，应当恳请最虔诚的皇帝下令，将整个非洲所有剩余的偶像完全清除；因为在许多沿海地区和不同的田产中，这一错误的邪恶仍然盛行：他们应命令将这些偶像清除，并下令将其庙宇（那些设在田间或偏僻之地、毫无装饰的）完全摧毁。\par

% structure: p0288 source=h3 target=subsection reason=relative-heading-rank; base=h2

\subsection{第五十九（希腊编号第六十三）条}

不得强迫圣职人员就其自身判决的审理公开作证\par

还应恳请他们恩准：若有人可能愿意按照赋予教会的宗徒法律在某个教会申辩其案件，并且若一方不满意圣职人员的裁决，则不得传唤担任审理者或在场的那位圣职人员出庭作证，也不得强迫任何依附于某位教会人员的人作证。\par

% structure: p0291 source=h3 target=subsection reason=relative-heading-rank; base=h2

\subsection{第六十（希腊编号第六十三）条}

关于异教庆典\par

还应恳请：由于在许多地方违背神圣训令举行庆典，这些庆典是由异教错误引入的，以致现在基督徒被迫由异教徒庆祝这些庆典；由此发生的情况是，在基督教皇帝的时代，似乎秘密地兴起了一种新的迫害；他们应命令禁止此类事情，并在刑罚之下禁止在城市和领地举行；尤其应当这样做，因为即使在最蒙福的殉道者的诞辰之日，在神圣的场所本身，他们也不惧犯下如此不义之事。因为在这些日子里，令人羞愧的是，他们在田野和空地举行最邪恶的跳跃，以至于妇女的尊严和无数因敬虔来到最圣日的女子的端庄受到放荡伤害的侵犯，以致几乎所有人都逃避对神圣宗教本身的接近。\par

% structure: p0294 source=h3 target=subsection reason=relative-heading-rank; base=h2

\subsection{第六十一（希腊编号第六十四）条}

关于演出，不得在主日或圣人庆节举行\par

此外，必须恳请：戏剧表演和其他演出的展示应从主日和基督教宗教的其他最神圣的日子中移除，尤其是在圣复活节八日庆期[即低主日]，人们聚集在竞技场多于在教堂；如果这些演出正好落在一个敬虔的日子，应将它们改到另一天，并且不得强迫任何基督徒观看这些演出，尤其是因为在执行违背天主训令的事情时，任何人都不得进行迫害，而（如理当如此）人应运用天主所赋予的自由意志。还特别应考虑合作者的危险，他们违背天主的训令，被极大的恐惧所迫去观看表演。\par

% structure: p0297 source=h3 target=subsection reason=relative-heading-rank; base=h2

\subsection{第六十二（希腊编号第六十五）条}

关于被定罪的神职人员\par

同时，应当请求他们俯允颁令：若有任何品级之圣职人员，因任何罪行而受主教判决，则其曾管辖之教会或任何人均不得为之辩护，违者将丧失钱财及职务。并且，命令不得以年龄或性别为借口而不予受理。\par

% structure: p0300 source=h3 target=subsection reason=relative-heading-rank; base=h2

\subsection{第六十三则（希腊第六十六则）}

关于成为基督徒的戏子\par

亦当请求他们：若有人愿从任何戏弄技艺（\OriginalTerm{戏弄技艺}{ludicra arte}）前来归依基督信仰之恩宠，并愿脱离那玷污，则任何人不得诱使或强迫他重返从事原先之事。\par

% structure: p0303 source=h3 target=subsection reason=relative-heading-rank; base=h2

\subsection{第六十四则（希腊第六十七则）}

关于在教会内举行解放奴隶仪式，须请求皇帝许可\par

关于在教会内公布解放奴隶之事，若我同道主教们在意大利各地亦行此礼，则遵循其程序将是我们信赖的标志。我们授予所遣特使全权，尽其所能，凡合乎信德，有益于教会状况及灵魂得救之事，我们将在主面前欣然接纳。凡此种种，若蒙圣座悦纳，请予明示，以使我能确认所提建议已获认可，并使这些诚心之举可自由接受我们一致的行动。全体主教说：圣座所命且所明智陈述者，皆合众心。\par

% structure: p0306 source=h3 target=subsection reason=relative-heading-rank; base=h2

\subsection{第六十五则（希腊第六十八则）}

关于被定罪的主教厄奎提乌斯\par

奥勒留主教说：我认为在使团中不应忽略厄奎提乌斯之事，此人早前已因罪行而受主教判决。若我方特使在那些地区巧遇此人，我兄弟当为教会的利益，见机行事，或在其能及之处，采取行动反对他。全体主教说：这项追究令我们极为满意，尤其厄奎提乌斯早已被定罪，他那无耻的躁动必须愈加在各处加以制止，以利于教会的善益与健康。他们签署：我，迦太基教会主教奥勒留，已同意此法令，读毕后签署姓名。同样，其他主教也都签署。\par

% structure: p0309 source=h3 target=subsection reason=relative-heading-rank; base=h2

\subsection{第六十六则（希腊第六十九则）}

应温和对待多纳特派人士\par

继而，在考虑了并处理了所有似乎有利于教会利益之事后，天主的圣神提示并劝勉我们，我们决定对那些前述人士——虽因无纪律的分裂而与主身体的合一隔绝——以温和与和平的方式行事，以便（尽我们所能）让所有那些被他们共融与团契之网所捕获的人，在非洲各省得知，他们如何因可悲的错误而执迷不悟。一如宗徒所说：\Quote{或许}当我们以温柔纠正他们时，\BibleQuote{天主会赐予他们悔改，得以认识真理，并摆脱魔鬼的罗网，因魔鬼已俘虏他们，随其意愿行事。}\BibleRef{弟后 2:25-26}\par

% structure: p0312 source=h3 target=subsection reason=relative-heading-rank; base=h2

\subsection{第六十七则（希腊第七十则）}

关于致函法官，请其记录多纳特派与马克西米安派之间所行之事\par

因此，我们认为应当由我们的公会议致函非洲的法官，由他们适宜地请求此事，即在此事上，他们应辅助共同的母亲——天主教会，使主教权威在各城得以巩固；亦即，凭其司法权力，并出于基督徒信德之勤勉，他们应调查并在公共案卷中记录各地点所发生之事——在这些地点，那些从他们中分裂出来的马克西米安派已获得了圣堂——以便众人对此有确切认知。\par

% structure: p0315 source=h3 target=subsection reason=relative-heading-rank; base=h2

\subsection{第六十八则（希腊第七十一则）}

多纳特派圣职人员准予作为圣职人员被接纳入天主教会\par

此外，我们认为应当致函我等同道主教弟兄，尤其是致宗座，即由我们前述可敬的弟兄与同僚阿纳斯塔修斯所主持的宗座，以便［希腊文\OriginalTerm{既然}{\GreekText{ἐπειδὴ}}，拉丁文\OriginalTerm{quo}{quo}］他知晓非洲极需教会的和平与繁荣：即那些曾是圣职人员的多纳特派人士，若出于善意渴望返回天主教合一，应按当地管理教会的主教之意愿与判断加以对待；若为基督徒和平之故，可保留其品级接纳，正如该分裂的先前时期已明显如此行。此事之例证，由多数乃至几乎所有此谬误曾产生的非洲教会所证实；并非要废除外国会议就此所召开之决议，而是该决议对如此自愿回归天主教之人仍保持效力，只要他们不导致合一的破裂。但对于那些在其居住地，藉明显赢得弟兄灵魂而促成或协助成全天主教合一之人，则不得援引外国会议所设反对其品级的法令，因为救恩不向任何人关闭；也就是说，多纳特派一方所祝圣之人，若经改正后愿回归天主教会，按照外国会议，不得按其品级接纳；但那些藉其接受建议回归天主教合一者，则为例外。\par

% structure: p0318 source=h3 target=subsection reason=relative-heading-rank; base=h2

\subsection{第六十九则（希腊第七十二则）}

为谋求和平，应向多纳特派派遣使团\par

此外，会议进一步决定，当这些事项完成后，应从我们中间派遣特使前往多纳图派（\OriginalTerm{多纳图派}{Donatists}）中那些被他们视为主教的人，或前往民众那里，为宣讲和平与合一，因为没有和平与合一，基督徒的救恩便无法获得；并且这些特使应引导所有人注意，他们对天主教会并无合理的反对理由。特别是，应通过市政文件（因其文书的权威性）向所有人显明，他们自己在对待自己内部的裂教者——马克西米安派（\OriginalTerm{马克西米安派}{Maximianists}）——一事上做了什么。因为在此事上，神圣的恩宠向他们显示——只要他们愿意留意——他们从教会合一中的分离，与他们如今所宣告的马克西米安派从他们中的分裂是同样不义的。然而，他们从数目中，将那些由他们全体会议的权威所谴责的人，连同其荣誉一起接纳回来，并接受了那些在被谴责和断绝期间所施行的洗礼。因此，让他们看看，当他们以愚蠢的心抵制散布于全世界的教会的和平时，他们从\Person{多纳图}（\Person{Donatus}）方面做了这些事，既不说自己因与那些为缔造和平而如此接纳的人共融而受玷污，却鄙视我们，即天主教会——它甚至建立在地极——，认为我们因与那些控告者无法赢到己方的人共融而被玷污。\par

% structure: p0321 source=h3 target=subsection reason=relative-heading-rank; base=h2

\subsection{第七十条（希腊编号第七十三条）}

哪些神职人员应与妻子禁欲\par

此外，由于有些神职人员被指控对自己的妻子不节制，会议决定：主教、长老和执事应根据既定的法规，甚至与自己的妻子禁欲；若他们不这样做，应被免去神职职务。但其余神职人员不应受此强制，而应遵循各教会在这方面的习惯。\par

% structure: p0324 source=h3 target=subsection reason=relative-heading-rank; base=h2

\subsection{第七十一条（希腊编号第七十四条）}

关于那些忽视自己教民的人\par

此外，会议决定：任何人不得离开其主教座堂（\OriginalTerm{主教座堂}{cathedral}）前往教区内所建的其他教堂，也不得因过度关心自己的事务而忽视对自己主教座堂的照顾和经常出席。\par

% structure: p0327 source=h3 target=subsection reason=relative-heading-rank; base=h2

\subsection{第七十二条（希腊编号第七十五条）}

关于在怀疑婴儿是否已受洗时为他们施洗\par

同样，会议决定：凡遇无可靠证人能毫无疑虑地证明婴儿已受洗，且婴儿本身因其幼龄而无法对圣事的施行作出回答时，所有此类婴儿应毫无疑虑地受洗，以免犹豫使他们失去圣事的洁净。这是由我们的弟兄、摩尔人（\OriginalTerm{摩尔人}{Moorish}）特使所敦促的，因为他们从野蛮人手中赎回许多这样的婴儿。\par

% structure: p0330 source=h3 target=subsection reason=relative-heading-rank; base=h2

\subsection{第七十三条（希腊编号第七十六条）}

应宣布复活节日期和会议日期\par

同样，会议决定：可敬的复活节日应通过格式信函（\OriginalTerm{格式信函}{formed letters}）的签署向所有人宣布；关于会议日期，亦应遵守同样的做法，即根据\Place{希玻}（\Place{Hippo}）会议的法令，定为九月朔日前第十日（\OriginalTerm{九月朔日前第十日}{X. Calends of September}），并应致函各省的首席主教（\OriginalTerm{首席主教}{primates}），使他们召集会议时不阻碍此日期。\par

% structure: p0333 source=h3 target=subsection reason=relative-heading-rank; base=h2

\subsection{第七十四条（希腊编号第七十七条）}

任何担任临时主教（\OriginalTerm{临时主教}{intercessor}）的人不得保留其担任临时主教的教座（\OriginalTerm{教座}{see}）\par

同样，会议规定：任何临时主教不得通过民众骚动和叛乱来保留他被任命为临时主教的那个教座；他应负责在一年内为他们提供一名主教；若他忽视此事，一年期满后，应任命另一位临时主教。\par

% structure: p0336 source=h3 target=subsection reason=relative-heading-rank; base=h2

\subsection{第七十五条（希腊编号第七十八条）}

关于向皇帝请求教会的保护人（\OriginalTerm{保护人}{defenders}）\par

鉴于穷人的苦难——他们的困扰使教会不断受到折磨——全体会议决定：应请求皇帝允许在主教监督下为穷人选出对抗富人势力的保护人。\par

% structure: p0339 source=h3 target=subsection reason=relative-heading-rank; base=h2

\subsection{第七十六条（希腊编号第七十九条）}

关于不出席会议的主教\par

同样，会议决定：每当应召集会议时，那些未因年迈、疾病或其他重大需要而受阻的主教应前来相聚，并应通知各省的首席主教（\OriginalTerm{首席主教}{primates}），以便从所有主教中组成两到三个小组，每个小组选出若干人，在会议日期前随时准备就绪；但若他们不能出席，应在传召书（\OriginalTerm{传召书}{tractory}）中写明理由，或若在传召书到达后突然发生某些紧急情况，除非他们向自己的首席主教报告其障碍，否则应满足于自己教会的共融。\par

% structure: p0342 source=h3 target=subsection reason=relative-heading-rank; base=h2

\subsection{第七十七条（希腊编号第八十条）}

关于\Person{克雷斯科尼乌斯}（\Person{Cresconius}）\par

关于\Place{王庄}（\Place{Villa Regis}）的\Person{克雷斯科尼乌斯}（\Person{Cresconius}），全体会议决定如下：应告知\Place{努米底亚}（\Place{Numidia}）的首席主教此事，以便他通过信函召叫上述\Person{克雷斯科尼乌斯}，使他在下一次非洲全体会议上不得拖延出席。若他藐视召叫而不来，则应让他知道应对他宣告判决。\par

% structure: p0345 source=h3 target=subsection reason=relative-heading-rank; base=h2

\subsection{第七十八条（希腊编号第八十一条）}

关于\Place{希玻-迪亚里图斯}（\Place{Hippo-Diarrhytus}）教会\par

会议进一步认为，既然希波-迪亚里图斯教会的匮乏不应再被忽视，且那些拒绝了埃基季乌斯可耻共融的人仍占有那里的教会，故应由本届会议派遣某些主教，即：\Person{雷吉努斯}、\Person{阿利比乌斯}、\Person{奥斯定}、\Person{马特尔努斯}、\Person{特亚修斯}、\Person{埃沃迪乌斯}、\Person{普拉西安}、\Person{乌尔班}、\Person{瓦莱里乌斯}、\Person{安比维乌斯}、\Person{福尔图纳图斯}、\Person{宽德武尔图斯}、\Person{霍诺拉图斯}、\Person{雅努阿里乌斯}、\Person{阿普图斯}、\Person{霍诺拉图斯}、\Person{安佩利乌斯}、\Person{维克托里安}、\Person{埃万格卢斯}和\Person{罗加提安}；当这些人聚集起来，并纠正了那些因有罪的固执而认为应等待同一埃基季乌斯逃走的人之后，应借众人投票为他们祝圣一位主教。但如果他们不愿考虑和平，则不应阻止为一位主教的祝圣而进行选举，以利于那长期匮乏的教会。\par

% structure: p0348 source=h3 target=subsection reason=relative-heading-rank; base=h2

\subsection{Canon 79. (Greek lxxxii.)}

\Emph{论那些未在一年内尽力使自己的案件得到审理的圣职人员}\par

会议进一步规定，凡被定罪并承认任何罪行的圣职人员，或因\OriginalTerm{受到宽恕的耻辱的他们}{eorum, quorum verecundiæ parcitur}，或因教会所受的侮辱，以及异端者和外邦人的傲慢夸耀，若他们偶然愿意出席自己的案件并声明自己无罪，应在被绝罚后一年内这样做；如果他们确实在一年内疏忽而未澄清自己的案件，此后将不再听取他们的声音。\par

% structure: p0351 source=h3 target=subsection reason=relative-heading-rank; base=h2

\subsection{Canon 80. (Greek lxxxiii.)}

\Emph{不允许任命非属自己管辖的隐修院的隐修士为隐修院院长，也不允许将他们祝圣为圣职人员}\par

同样，会议认为，若任何主教愿意将从非其管辖的隐修院接收的隐修士擢升为圣职人员，或任命他为所属隐修院的院长，则如此行事的主教应与他人断绝共融，只满足于与自己教区子民的共融；但那隐修士不得继续担任圣职人员或院长。\par

% structure: p0354 source=h3 target=subsection reason=relative-heading-rank; base=h2

\subsection{Canon 81. (Greek lxxxiv.)}

\Emph{论指定异端者或外教人为继承人的主教}\par

同样，会议规定，若有任何主教偏爱与他不具血缘关系的陌生人、或他的异端亲属、或外教人作为他的继承人，而非教会，则他即使在死后也应受\OriginalTerm{绝罚}{anathema}，其名字绝不可在天主司铎的名录中被诵读。若他未立遗嘱而死，也不得被宽恕，因为身为主教，他本不应延迟作出适合其身份的财产处置。\par

% structure: p0357 source=h3 target=subsection reason=relative-heading-rank; base=h2

\subsection{Canon 82. (Greek lxxxv.)}

\Emph{论解放奴隶}\par

同样，会议认为应恳请皇帝关于在教会中宣布解放奴隶之事。\par

% structure: p0360 source=h3 target=subsection reason=relative-heading-rank; base=h2

\subsection{Canon 83. (Greek lxxxvi.)}

\Emph{论虚假的殉道者纪念地}\par

同样，会议认为，那些在田野和路旁作为殉道者纪念地随处设立的祭台，若无法证明其中安放了任何殉道者的遗体或圣髑，应由管理该地区的主教予以拆除——如果可行的话。但若因可能引起民众骚动而无法做到，则应告诫民众不要前往这些地方，且信仰正确的人不应受任何地方迷信的束缚。除非发现有遗体或某些圣髑，并且根据传统和可靠权威最初被宣称为其住所、产业或受难地点，否则绝不可接受任何殉道者纪念地。因某些人的梦或空洞的\Emph{准}启示而在任何地方设立的祭台，应一概予以否定。\par

% structure: p0363 source=h3 target=subsection reason=relative-heading-rank; base=h2

\subsection{Canon 84. (Greek lxxxvii.)}

\Emph{论清除偶像的残余}\par

同样，会议认为应恳请最荣耀的皇帝们，将偶像崇拜的残余不仅在雕像上，而且在任何地方、小树林或树木上，全部予以清除。\par

% structure: p0366 source=h3 target=subsection reason=relative-heading-rank; base=h2

\subsection{Canon 85. (Greek lxxxviii.)}

\Emph{论必要时由迦太基主教以全体主教的名义撰写并签署信函}\par

全体主教说：若需以会议名义撰写任何信函，会议认为，主持此教座的可敬主教应惠予代表全体口述并签署之，其中包括那些要派往非洲各省的主教代表关于多纳徒派事宜的信函；同时认为，交给他们的信函应包含委托书的要旨，他们不得逾越。他们签署如下：\Quote{我，迦太基教会的主教\Person{奥勒留}，已同意此法令，阅读后已签署。}同样，其余所有主教也签署了。\par

……\par

本次会议确认了先前的法令。\par

在最荣耀的皇帝\Person{阿卡狄乌斯}和\Person{霍诺留}（奥古斯都）的第五次执政期间，九月\OriginalTerm{第六朔日}{VI Calends of September}，在\Place{米莱维斯}城的大殿会议室中，当\Place{迦太基}主教\Person{奥勒留}在全体会议上就座，执事侍立时，主教\Person{奥勒留}说：\Quote{既然神圣教会为同一身体，且所有肢体的头只有一个，因天主的许可和对我们软弱的助佑，我们受兄弟之爱的邀请来到此教会。因此，我恳求你们的爱德，相信我们到你们这里来既非多余，也不为众人所不悦；我们众人的同意将表明，我们同意已在\Place{希波}会议中确认或后来在\Place{迦太基}更大主教会议上所规定的法令，这些法令现在应按顺序向我们宣读。最后，若他们知道你们不仅因同意这些文件，也因你们的署名而公布了我们在先前会议中合法规定的事项，则你们圣洁的同意将比光更显明。}\par

努米底亚首席教座的主教\Person{桑提普斯}说：\Quote{我相信，凡令众弟兄和法规所喜悦且亲手确认的事，我们通过签署姓名也表示赞同，并已用我们的签名确认。}\par

\Person{尼切提乌斯}，毛里塔尼亚·西提芬西斯首座之主教说：所宣读之法令，既不乏理据，且为众人所赞同，我之卑微亦悦纳之，并将以署名确认。\par

% structure: p0374 source=h3 target=subsection reason=relative-heading-rank; base=h2

\subsection{第八十六条教规（希腊第八十九条）}

\Emph{关于主教之秩序：较晚受圣秩者不得僭越于较早受圣秩者之前。}\par

\Person{瓦伦廷}主教说：若蒙贵方忍耐允准，我遵循往昔在迦太基教会中所行之事，且这些事因诸位兄弟之署名而得以确认，显赫昭著，我声明我们决意予以保存。因我等知晓，教会纪律始终未受侵犯：因此，任何兄弟不得胆敢将自己置于较己早受圣秩者之前；然藉爱德之职事，此礼一向向较早受圣秩者显呈，且较晚受圣秩者常应欣然接纳。请尊位颁令，以诸君之表明巩固此秩序。\Person{奥勒留}主教说：若非存在某些轻率之心，诱使我们制定此类法规，本不应重述此事；然此乃我兄弟、同为主教者所言及之共同缘由，即我们每位应承认天主所定之秩序，较晚者应服从较早者，且当后者未经咨询时，不得擅自作为。是以我言，今思及之，凡自以为可僭越于较早受圣秩者之前者，应由全体会议加以适当约束。\Person{克桑提普斯}，努米底亚首座主教说：出席之全体兄弟已闻我兄弟、同为主教者\Person{奥勒留}所言，吾等作何答复？\Person{达提安}主教说：祖先所定之法令应藉我等同意而强化，如此迦太基教会在往昔会议中所行之事，经我等全体一致确认，应持守不渝。全体主教说：此秩序已由我父辈及祖先保存，我等亦将赖天主助佑保存之，同时保持努米底亚及毛里塔尼亚首席权之权利不受侵犯。\par

\Emph{关于努米底亚之档案及名册。}\par

此外，凡在此次会议中署名之全体主教认为：努米底亚之名册及档案应置于首座及都会，即君士坦丁纳。\par

\Signature{约翰逊。}\par

\EditorialNote{由此教规可见，非洲之首席权具流动性，属于该省资历最老之主教。若首席权固定于某特定城市主教，如他处然，则应有为该主教设定之保留或例外，一如西班牙布拉加主教会议第二十四条教规所定，该教规命令全体主教依其资历居位，但保留都会主教之权利。迦太基主教不在本条教规范围内；因显然他享有附属于其主教座之优先权，实则类似宗主教。努米底亚及毛里塔尼亚被特别提及之原由，乃因彼处曾就该问题发生若干争议。}\par

% structure: p0381 source=h3 target=subsection reason=relative-heading-rank; base=h2

\subsection{第八十七条教规（希腊第九十条）}

\Emph{关于主教\Person{克沃德武尔图斯}。}\par

关于森图里亚之\Person{克沃德武尔图斯}案，全体主教议决：在其案件终结前，任何人不得与其共融；因其控告者曾欲将案件提交本会议，经询问是否愿与该主教在主教们面前受审，初时表示愿意，后一日却答称不愿，遂离去。在此情形下，于知其案件结果之前，即剥夺其主教职位，此事无任何基督徒可视为公允之举。\par

% structure: p0384 source=h3 target=subsection reason=relative-heading-rank; base=h2

\subsection{第八十八条教规（希腊第九十一条）}

\Emph{关于主教\Person{马克西米安}。}\par

至于瓦盖之\Person{马克西米安}案，会议认为宜致函其本人及其教民，命其卸任主教之职，并请他们另求指定他人。\par

% structure: p0387 source=h3 target=subsection reason=relative-heading-rank; base=h2

\subsection{第八十九条教规（希腊第九十二条）}

\Emph{凡受圣秩之主教，应由其授秩者领取载明日期及执政官名衔之信函。}\par

此外，会议进一步议决：此后在非洲各省内由主教祝圣者，应由其授秩者领取亲手写成之信函，内载执政官名衔及日期，以免就何人先受圣秩、何人在后而起争议。\par

% structure: p0390 source=h3 target=subsection reason=relative-heading-rank; base=h2

\subsection{第九十条教规（希腊第九十三条）}

\Emph{凡曾在教会诵读一次者，不得由他人提升。}\par

同样，会议议决：凡在教会中曾诵读一次者，不得由另一教会接纳进入圣职（\OriginalTerm{圣职}{clericatum}）。\par

他们署名如下：我，迦太基教会主教\Person{奥勒留}，已同意此法令，阅读后签字。其余主教亦同样签名。\par

\Emph{本次会议陈述了被派往海外为特使之主教所行之事。}\par

在最显赫者——最荣耀之皇帝\Person{德奥多西·奥古斯都}与\Person{鲁莫里杜斯}任执政官之年份，八月二十四日，在迦太基第二区之巴西利卡中，当\Person{奥勒留}主教在全体会议上就座，执事侍立时，\Person{奥勒留}主教说：尊敬之诸兄弟，缘于形势所迫，我虽卑微，却引导你们集会于此。因前些时日，如诸尊位所忆，我们举行会议期间，曾派遣我兄弟为特使往海外地区。此时诸尊位聚会，理应听其陈述现已完成之特使使命所经之历程。尽管昨日当我们就此事开会时，除教会事务外，已对其所行之事予以一定持续之关注，然今日仍应以教会行动确认昨日之讨论。\par

\Emph{关于未出席本次会议之非洲各省主教。}\par

事物的正确秩序要求我们首先应询问我们的兄弟及同僚主教，他们本应前来参加这次公会议，或是来自\Place{Byzacena}，或至少来自\Place{Mauritania}，如同他们曾决定将出席本次会议。当\Place{Byzacena}省的主教\Person{Philologius}、\Person{Geta}、\Person{Venustianus}和\Person{Felician}呈递并宣读了他们的使节信函，而\Place{Mauritania Sitiphensis}省的使节\Person{Lucian}和\Person{Silvanus}也做了同样的事之后，主教\Person{Aurelius}说：\Quote{请将这些信件的原文载入会议记录。}\par

\Emph{论\Place{Byzacene}主教}\par

主教\Person{Numidius}说：\Quote{我们注意到我们的兄弟及同僚，即\Place{Byzacena}省和\Place{Mauritania Sitiphensis}省的主教，已派遣使节前来公会议；现在我们询问，\Place{Numidia}的使节是否已经到来，或至少\Place{Tripoli}省或\Place{Mauritania-Cæsariensis}省的使节是否到来。}\par

\Emph{论\Place{Mauritania Sitiphensis}的主教}\par

主教\Person{Lucian}和\Person{Silvanus}，\Place{Mauritania Sitiphensis}省的使节，说：\Quote{通谕函件送达我们\Place{Cæsarian}的兄弟们太迟，否则他们本会到场；而且他们一定会来，我们相信他们的心意，即凡此次公会议所决定的，他们无疑将赞同。}\par

\Emph{论\Place{Numidia}的主教}\par

\Place{Tagaste}教会的主教\Person{Alypius}说：\Quote{我和圣洁的弟兄\Person{Augustine}与\Person{Possidius}来自\Place{Numidia}，但无法从\Place{Numidia}派出使节，因为新兵的骚乱使得主教们要么无法前来，要么忙于各自教区必需的事务。因为在我将尊驾的通谕函件呈交圣洁的长老\Person{Xantippus}之后，他认为在这件事上应当指定一次公会议，并派遣带有指示的代表团前往；但当我后来在信件中将我刚才提到的招募新兵的阻碍报告给他时，他以自己的谕令免除了他们。}主教\Person{Aurelius}说：\Quote{毫无疑问，上述\Place{Numidia}的弟兄和主教们在收到公会议记录后，将表示同意，并努力执行任何被采纳的决定。因此，有必要通过本座的关怀，将我们所决定的事通知他们。}\par

\Emph{论\Place{Tripoli}的主教}\par

\Quote{以下是我能获知的关于\Place{Tripoli}弟兄们的事：他们指定了我们的兄弟\Person{Dulcicius}为使节；但由于他未能前来，一些来自该省的我们的儿子们声称，上述之人已经登船，并且人们认为他的抵达因风暴而延迟；尽管如此，关于这些事，若尊驾愿意，将保留这一形式，即把公会议的决议送交他们。}所有主教都说：\Quote{尊驾所决定的令我们所有人都满意。}\par

% structure: p0406 source=h3 target=subsection reason=relative-heading-rank; base=h2

\subsection{Canon 91. (Greek xciv.)}

\Emph{论与多纳徒派举行会议}\par

主教\Person{Aurelius}说：\Quote{在尊驾的处理中出现的事，我认为应当由教会法令加以确认。因为你们所有人的声明表明，我们每一人都应在自己的城市中召集多纳徒派的首领，或是单独，或是与一位邻近的主教一起，以便同样在各城市和地方，由该地方的官吏或长者召集他们举行会议。若各位认为合适，便将此制成谕令。}所有主教说：\Quote{我们都认为合适，并且我们都已用签名确认了这一点。此外，我们恳请尊驾签署拟由公会议送交法官的信函。}主教\Person{Aurelius}说：\Quote{若尊驾认为合适，请宣读召集他们的格式，以便我们都遵循相同的诉讼程序。}所有主教说：\Quote{请宣读。}文书\Person{Lætus}宣读。\par

% structure: p0409 source=h3 target=subsection reason=relative-heading-rank; base=h2

\subsection{Canon 92. (Greek xcv.)}

\Emph{召集多纳徒派的格式}\par

那位教会的主教说：\Quote{我们请求尊驾宣读凭借那最尊荣宗座的权威所获得的内容，并命令将其载入记录并执行。}当\OriginalTerm{命令}{jussio}被宣读并载入记录后，天主教会的主教说：\Quote{请垂听通过尊驾送交多纳徒派的训令，并将其载入记录，送至他们那里，并告知我们他们在记录中的答复。}\Quote{我们受公会议的权威派遣，召集你们前来，渴望因你们的矫正而喜乐，谨记主的仁爱，祂曾说过：\BibleQuote{缔造和平的人是有福的，因为他们要称为天主的子女。}；再者，祂借着先知告诫那些说他们不是我们兄弟的人，我们应当说：\InnerQuote{你们是我们的弟兄。}因此，你们不应轻视这出自爱德的和平劝谕，以至于若你们认为我们拥有任何真理的部分，你们不犹豫地说出来：也就是说，当你们的会议聚集时，你们从你们中委派一些人，将陈述你们案件的事托付给他们；好使我们也能这样做：从我们的会议中委派一些人，与你们委派的人一起，在确定的地点和时间，和平地讨论任何使你们的共融与我们分离的问题；最终，借着我们主天主的助佑，旧日的错误得以终结，免得因人的仇视，软弱的灵魂和愚昧的人因亵圣的分裂而丧亡。但如果你们以兄弟的精神接受此提议，真理将轻易显明；但如果你们不愿意这样做，你们的不信任将容易被知道。}当此被宣读后，所有主教说：\Quote{这使我们十分满意，就这样行。}他们签署：\Quote{我，\Person{Aurelius}，\Place{Carthage}教会的主教，同意此法令，在阅读后签署了它。}同样，其余的主教也签署了。\par

\Emph{本次会议向君主们派遣了反对多纳徒派的使节}\par

最光荣的皇帝\Person{Honorius Augustus}，在担任第六次执政官时，于七月一日，在\Place{Carthage}的第二区大殿。在这次会议中，\Person{Theasius}和\Person{Euodius}接受了反对多纳徒派的使节任务。这次会议中插入了随后的一份\OriginalTerm{劝谕书}{commonitorium}。\par

% structure: p0414 source=h3 target=subsection reason=relative-heading-rank; base=h2

\subsection{第九十三则（希腊第九十六则）}

\Emph{使节所领受反对多纳徒派之训谕书的性质}\par

我们为兄弟\Person{特阿西乌斯}和\Person{埃沃迪乌斯}所拟的训谕书，他们是由\Place{迦太基}会议派往最光荣最虔诚的君王们那里作使节的。当赖主的助佑，他们觐见最虔诚的君王时，应怀着充分的确信，按照前年会议指示，宣告多纳徒派的主教们曾被市镇当局催促聚会，以便如果他们真愿实践他们所宣称的，他们便可从他们中选出适当的人，以基督般的温顺与我们进行和平会议，并毫不犹豫地坦诚宣告他们所持为真理的；如此，从这次会议中，天主教会那早已显明的真诚，甚至能被那些因无知或固执而反对它的人所察觉。但因缺乏信心，他们几乎不敢回答。而实际上，因为我们向他们履行了主教与和平缔造者所当行的职务，他们对于真理却无话可说，他们便诉诸无理性的暴力行为，背信弃义地压迫了许多主教和神职人员，更不用说平信徒了。有些教堂他们竟实际侵占了，并且试图攻击其他教堂。\par

现在，陛下的仁慈恩惠应当采取措施，使天主教会得以在其庇护下得到保护，因为天主教会曾在她的怀抱中孕育他们作为基督的敬拜者，并以信德的刚强之食养育他们；以免暴力之徒利用宗教狂热之时，以恐惧压倒那些他们无法以辩论诱惑的软弱人民。众所周知，锡康塞利翁派的可憎之群（\OriginalTerm{可憎之群}{detestabilis manus}）曾屡遭法律禁止，并受其陛下最虔诚的诸位先皇的多项谕令谴责。针对这等人的疯狂，请求世俗（希腊文\OriginalTerm{神圣}{\GreekText{θείας}}）的保护既非异常，亦不违圣经，因为正如正典\BibleBook{宗徒}{大事录}所记载的，\Person{保禄}宗徒曾藉军事帮助避开了某些不法之徒的阴谋。因此，现在我们请求，毫无掩饰地向各城市的驻军（拉丁文\OriginalTerm{驻军}{ordinum}）以及各地郊野地产的所有者提供保护。同时，有必要请求陛下颁布命令，重申由他们值得纪念的父亲\Person{狄奥多西}所颁布的法律：对异端者的祝圣者与被祝圣者均处以十磅黄金的罚款，并针对在其宅邸举行秘密聚会的地主；如此，当天主教徒因受其诡计激怒而提起诉讼时，该法律可对此类人等生效；以便至少因恐惧，他们可以停止制造分裂和异端者的邪恶行为，因为他们拒绝因永罚的思考而被洁净和纠正。\par

也请求陛下立即更新法律，剥夺异端者接收或遗赠礼物或遗嘱的能力，以便那些因固执的疯狂而盲目、决意继续留在多纳徒派错误中的人，所有给予或接收的权利均被剥夺。\par

至于那些因合一与和平的考虑而愿意改正自己的人，应准许他们接收其遗产，尽管有法律禁止，即使该礼物或遗嘱的遗赠是在他们尚处于异端错误中时所作；当然，那些因诉讼压力而寻求进入天主教会的人除外；因为很可能，后一类人渴望大公教会合一并非因畏惧天主的审判，而是出于贪恋世俗的利益。\par

为推进这一切，需要每个省份的官员（\OriginalTerm{权力}{Potestatum}）的帮助。至于其他事项，凡他们认为有利于教会之事的，我们决定授予代表团全权执行并落实。此外，我们全体决定，应由我们的会议向最光荣的皇帝们和最卓越的尊贵们发出信函，使他们确信我们全体同意代表团应由我们派往其至福宫廷。\par

让我们全体在此信函上签名是一件非常缓慢的事，而为了不令他们为我们每个人的署名所累，我们希望你，\Person{奥勒留}兄弟，请你的爱德代我们全体签名。对此全体同意。\par

我，\Person{奥勒留}，\Place{迦太基}教会主教，同意此项法令并已签名。其他所有主教也都签名。\par

也应向法官们发出信函，请求在主允许代表们回到我们这里之前，他们通过城市士兵和各处天主教会农场的管理者提供保护。还应补充关于不诚实的\Person{厄奎提乌斯}的事，他因主张司铎的权利而显出其不诚实，应按照皇帝法令予以拒绝，不得进入\Place{希波}教区。还应向\Place{罗马}教会主教发出信函，推荐代表团，并向可能身在皇帝身边的其他主教发出信函。众人对此表示赞同。\par

同样，我，\Person{奥勒留}，\Place{迦太基}教会主教，同意此项法令，并阅读后已署名。\par

其他所有主教同样签名。\par

% structure: p0426 source=h3 target=subsection reason=relative-heading-rank; base=h2

\subsection{第九十四则（希腊第九十七则）}

\Emph{章节概要}\par

决定从各省向米佐尼姆（\Place{Mizoneum}）派遣自由代表团。命令派遣使节并发出信函，以指导自由使团：由于合一仅在迦太基实现，也应致函法官，命其在其他省份和城市继续进行合一工作，并且迦太基教会关于排除多纳徒派（\OriginalTerm{多纳徒派}{Donatists}）而对整个阿非利加的感恩之情应连同主教们的信函一并呈送宫廷（\OriginalTerm{宫廷}{ad Comitatum}）。\par

宣读了教宗依诺增爵（\Person{Innocent}）的信函：主教们不宜轻易将案件携往海外，此事本身也得到主教们自己判决的确认；为感谢并排除多纳徒派（\OriginalTerm{多纳徒派}{Donatists}），应派遣迦太基教会的两名圣职人员前往宫廷。\par

% structure: p0430 source=h3 target=subsection reason=relative-heading-rank; base=h2

\subsection{第九十五条（希腊第九十八条）}

\Emph{仅当必要时才召开普世会议}\par

以为善，不再年年疲累弟兄们；但每当共同事务，即整个阿非加的事务要求时，各地应在此事上致函该教座，以便在因需要而适宜召开会议的省份召集教务会议；至于不涉及共同利益之案件，应在各自省份内审理。\par

本条规则实质上撤销了先前会议所订第十八条规则中关于年度教务会议的条款。\par

% structure: p0434 source=h3 target=subsection reason=relative-heading-rank; base=h2

\subsection{第九十六条（希腊第九十九条）}

\Emph{对已选定的法官不得上诉}\par

若有人提起上诉，上诉人应选择法官，被上诉人也一同选择；他们的判决不得再上诉。\par

\Emph{关于各省的代表}\par

当各省的全体代表齐聚时，他们受到最友善的接待，即努米底亚人、比扎塞纳人、斯提芬斯摩尔人，以及凯撒利人和的黎波里人的代表。\par

\Emph{关于教会执行人}\par

此外，认为宜请求任命五名执行人，负责处理教会所需的一切事务，他们将被分派到各省。\par

% structure: p0441 source=h3 target=subsection reason=relative-heading-rank; base=h2

\subsection{第九十七条（希腊第一百条）}

\Emph{请求皇帝在教会诉讼中提供辩护人的保护}\par

认为宜由即将出发的使节味增爵（\Person{Vincent}）和福尔图纳提安（\Person{Fortunatian}）以全体省份的名义请求最荣耀的皇帝们授予设立司法辩护人（\OriginalTerm{司法辩护人}{scholastic defensores}）的权力，由其负责处理此类事务：如此，作为省司铎（\OriginalTerm{司铎}{priests}），他们获得教会事务中辩护人的权能，每当必要之时，可进入法官的私人居所，以反驳对方所强调的内容或进行必要的解释。\par

\Emph{毛里塔尼亚主教们反对普里莫苏斯（\Person{Primosus}）的抗议。}\par

显然，凯撒利毛里塔尼亚的主教们在其书面文件中证明，蒂加嫩西斯城（\Place{Thiganensian}）的首领曾传唤普里莫苏斯（\Person{Primosus}），命其依照帝国法令出席全体会议，并在他被依法搜求时，普里莫苏斯未被找到，至少执事们如此报告。但由于同一毛里塔尼亚主教们请求由全体教务会议致函可敬的弟兄、年迈的依诺增爵（\Person{Innocent}），认为宜发送此信，使他得知普里莫苏斯已在教务会议中被传唤而根本未找到。\par

% structure: p0446 source=h3 target=subsection reason=relative-heading-rank; base=h2

\subsection{第九十八条（希腊第一百零二条）}

\Emph{关于从未有过主教的民众}\par

以为善，凡从未有过自己主教的民众，绝不得接受主教，除非经各省全体会议及首席主教（\OriginalTerm{首席主教}{primates}）的决议，并得到该教会所属教区主教的同意。\par

% structure: p0449 source=h3 target=subsection reason=relative-heading-rank; base=h2

\subsection{第九十九条（希腊第一百零三条）}

\Emph{关于从多纳徒派（\OriginalTerm{多纳徒派}{Donatists}）回归的民众或教区}\par

凡从多纳徒派（\OriginalTerm{多纳徒派}{Donatists}）回归的团体已有主教者，无疑可继续拥有他们，即使没有教务会议的决议亦可；但若某团体曾有一位主教，当他去世后，该团体不再愿意拥有自己的主教，而愿归属另一主教之教区，此请求不应被拒绝。此外，凡在关于合一的帝国法律颁布之前已将所属民众带回天主教会的主教，应获准继续治理他们；但从该合一法律颁布之时起，所有教会及其教区，以及若有任何教会文书或属于其权利之物，无论其原属于仍为异端者还是后归正天主教会者，均应属于该地天主教会主教。凡在此法律之后有任何侵占行为者，应归还所侵占的产业。\par

% structure: p0452 source=h3 target=subsection reason=relative-heading-rank; base=h2

\subsection{第一百条（希腊第一百零四条）}

\Emph{关于毛伦提乌斯（\Person{Maurentius}）主教的建议}\par

[赫费勒（\Person{Hefele}）说：\Quote{本条规则的文本残缺甚多，极难理解。} 他给出的概述是：\Quote{会议为毛伦提乌斯（\Person{Maurentius}）主教一案指定了法官。}（赫费勒，第二卷，第443页。）]\par

约翰逊（\Person{Johnson}）如此概括并翻译：毛伦提乌斯主教有一项对他不利的指控已提交会议，他请求听证；但指控者在指定日期经执事三次传唤后仍未到庭。此案被提交给长者克桑提普斯（\Person{Xantippus}）、奥斯定（\Person{Augustinus}）以及由会议召集的其他五人，指控者应凑足十二人。\par

% structure: p0456 source=h3 target=subsection reason=relative-heading-rank; base=h2

\subsection{第一百零一条（希腊第一百零四条\OriginalTerm{bis}{bis}）}

\Emph{关于调节罗马教会与亚历山大里亚教会之间的和平}\par

以为善，应致函圣教宗依诺增爵（\Person{Innocent}），论及罗马教会与亚历山大里亚教会之间的不和，以便各教会彼此保持和平，正如主所命令的。\par

% structure: p0459 source=h3 target=subsection reason=relative-heading-rank; base=h2

\subsection{第一百零二条（希腊第一百零五条）}

\Emph{论休妻或休夫而保持独身者}\par

会议认为，照福音与纪律，被休的丈夫或妻子不应另娶或另嫁，而应保持现状，或彼此和好；若他们蔑视此律，则应强制他们做补赎。关于此案，我们必须请求颁布一条帝国法律。\par

% structure: p0462 source=h3 target=subsection reason=relative-heading-rank; base=h2

\subsection{第一〇三条（希腊第一〇六条）}

\Emph{关于在祭台上应念的祷文}\par

同样也认为，凡在主教会议中已获批准的祷词，包括序文、推荐词和按手礼，均应由众人使用；绝不可使用任何违反信德的祷词，而只应念诵那些由博学者所收集的。\par

% structure: p0465 source=h3 target=subsection reason=relative-heading-rank; base=h2

\subsection{第一〇四条（希腊第一〇七条）}

\Emph{关于那些请求皇帝让世俗法官审理自己案件的人}\par

会议认为，凡向皇帝请求让世俗法官审理自己事务者，应剥夺其职务；然而，若他请求皇帝给予主教审理，则不应反对。\par

% structure: p0468 source=h3 target=subsection reason=relative-heading-rank; base=h2

\subsection{第一〇五条（希腊第一〇八条）}

\Emph{关于不在非洲共融而欲渡海者}\par

凡不在非洲共融，而渡海去共融者，应开除其神职。\par

% structure: p0471 source=h3 target=subsection reason=relative-heading-rank; base=h2

\subsection{第一〇六条（希腊第一〇九条）}

\Emph{那些将去宫廷诉讼者应注意通知迦太基主教或罗马主教。}\par

会议认为，凡愿去宫廷者，应按照送交罗马教会的形式递交通知，以便从那里获得一份正式信函以转往宫廷。但若只领得一封致罗马的正式信函，却未提及自己需要去宫廷，便径往该处，则应与教会断绝共融。然而，若在罗马时突然出现去宫廷的需要，则应将此需要陈述于罗马主教，并携带该罗马主教的一份批复。至于由首主教及某些主教发给其本籍神职人员的正式信函，应注明复活节日期；但若该年复活节日期尚不确定，则应注明上一年的复活节日，正如公共文件通常以执政官年号记日一样。\par

此外会议认为，由本届光荣会议派出的代表应向最光荣的君王们请求他们认为有益于反对多纳特派与异教徒及其迷信的一切事宜。\par

所有主教也认为，所有会议信函应由您一人签署。于是他们签署如下：我，迦太基主教奥勒留，已同意此法令，经阅读后，现签署我的名字。同样，其余主教亦签署。\par

% structure: p0476 source=h3 target=subsection reason=relative-heading-rank; base=h2

\subsection{（希腊第一一〇条）}

\Emph{反对异教徒和异端者的会议}\par

在最显赫的巴苏斯和菲利普担任执政官之年，七月十六日，于迦太基修复大殿的议事厅。此次会议中，福图纳提安主教第二次被任命为反对异教徒和异端者的代表。\par

\Emph{另一次反对异教徒和异端者之会议}\par

在最显赫的巴苏斯和菲利普担任执政官之年，十月十三日，于迦太基修复大殿的议事厅。此次会议中，雷斯提图图斯和弗洛伦提乌斯两位主教接受了反对异教徒和异端者的使命，当时塞维鲁和马加略被杀，因他们之故，主教埃沃迪乌斯、泰阿修斯和维克托被处死。\par

% structure: p0481 source=h3 target=subsection reason=relative-heading-rank; base=h2

\subsection{第一〇七条（希腊第一一〇条续）}

\Emph{关于主教审理案件之会议}\par

在最光荣的皇帝奥诺留第七次与狄奥多西第三次任执政官之年，七月十六日，于迦太基第二区大殿举行了一次会议。此次会议认为，任何主教均不得单独主张审理案件的权力。此会议的法令我未在此写下，因为它只是省区性的，而非普世性的。\par

% structure: p0484 source=h3 target=subsection reason=relative-heading-rank; base=h2

\subsection{（希腊第一一一条）}

\Emph{反对多纳特派之会议}\par

在最光荣的皇帝奥诺留第八次与狄奥多西第四次任执政官之后，七月十八日，于迦太基第二区大殿。此次会议中，弗洛伦提乌斯、波西迪乌斯、普雷西迪乌斯和贝内纳图斯四位主教接受了反对多纳特派的使命，当时颁布了一项法律，允许任何人按其意愿实践基督教敬拜。\par

% structure: p0487 source=h3 target=subsection reason=relative-heading-rank; base=h2

\subsection{第一〇八条（希腊第一一二条）}

\Emph{反对佩拉纠和塞莱斯提乌斯异端之会议}\par

在最光荣的皇帝、至尊奥古斯都奥诺留第十二次与狄奥多西第八次任执政官之年，五月一日，于迦太基福斯图斯大殿的议事厅。在执事们侍立之时，奥勒留主教主持了整个会议。全体聚集于迦太基教会神圣会议中的主教，其姓名和签署均已列明，一致同意定义如下：\par

% structure: p0490 source=h3 target=subsection reason=relative-heading-rank; base=h2

\subsection{第一〇九条（希腊第一一二条续）}

\Emph{亚当并非由天主造成死亡所辖}\par

若有人说第一人亚当被造为可死的，以致无论他犯罪与否，他都会在身体上死亡——即离开身体，不是因他的罪所应得，而是出于自然必然性，则应受绝罚。\par

% structure: p0493 source=h3 target=subsection reason=relative-heading-rank; base=h2

\subsection{第一一〇条（希腊第一一二条之二）}

\Emph{婴孩受洗是为赦免罪过}\par

同样会议认为，若有人否认新从母腹出的婴孩应当受洗，或说洗礼虽是为赦免罪过，但婴孩从亚当身上并未继承任何原罪（需由重生之洗除去），从而得出结论说，在他们身上，为赦免罪过而施行的洗礼形式应被理解为虚假而非真实，则应受绝罚。\par

因为宗徒所说的\BibleQuote{罪恶藉着一人进入了世界，死亡藉着罪恶也进入了世界，这样死亡就传到了众人，因为众人都犯了罪}，若不这样理解，就无法明白。普世教会始终是这样理解的。因为根据这\OriginalTerm{信德的规诫}{regulam fidei}，连那些自己尚未犯任何罪的婴儿，也真正为了赦罪而受洗，以便藉着再生洗净他们因生育而来的东西。\par

据\Person{苏里乌斯}记载，在一份非常古老的抄本中此处有如下内容。这不见于希腊文本，也不见于\Person{狄奥尼修}。\Person{布伦斯}将其置于脚注中。\par

[同样，议定：若有人说主的这句话：\BibleQuote{在我父的家里有许多住处}应理解为在天国里将有一个中间的地方，或某个地方，在那里未受洗而离开此生的婴孩过得快乐，而没有圣洗（没有它就不能进入天国，即永生）的人，则此人应受绝罚。因为主曾说：\BibleQuote{人除非由水和圣神重生，不能进入天国}，哪个天主教徒会怀疑那没有堪当成为基督同继承者的人将变成魔鬼的同伙？因为凡不能得到右边的人，无疑将得到左边的份。]\par

% structure: p0499 source=h3 target=subsection reason=relative-heading-rank; base=h2

\subsection{第一一一条（希腊编号一一三）}

\Emph{天主的恩宠不仅赦罪，还提供帮助使我们不再犯罪}\par

同样，议定：凡说使人藉着我们的主耶稣基督而成义的天主恩宠，仅足以赦免已往的罪，而不能助人避免将来犯罪的人，应受绝罚。\par

% structure: p0502 source=h3 target=subsection reason=relative-heading-rank; base=h2

\subsection{第一一二条（希腊编号一一三续）}

\Emph{基督的恩宠不仅赐给我们认识本分的知识，还激发我们渴望能够完成我们所知道的}\par

同样，凡说藉着我们的主耶稣基督而来的同一恩宠，仅藉着启示和开启我们的心智以理解诫命来帮助我们不去犯罪，使我们知道该寻求什么、该避免什么，并且使我们乐意如此行，但并未藉着它获得帮助以能实行我们所知当行之事的人，应受绝罚。因为当宗徒说：\BibleQuote{知识使人自高自大，惟有爱德能建树人}时，我们若相信我们拥有了那使我们自高自大的恩宠，却没有那建树人的恩宠，那真是可耻，因为两者都是天主的恩赐：既知道该做什么，又乐意去做；这样，知识就不能在我们自高自大时爱德却建树我们。正如圣经论及天主说：\BibleQuote{祂教导人知识}，同样也写着：\BibleQuote{爱德出于天主}。\par

% structure: p0505 source=h3 target=subsection reason=relative-heading-rank; base=h2

\subsection{第一一三条（希腊编号一一四）}

\Emph{没有天主的恩宠，我们不能行任何善事}\par

议定：凡说成义的恩宠仅赐给我们，使我们能藉着自由意志更容易履行所命之事；好像即使不赐恩宠，虽不容易，但我们仍能没有恩宠而完成天主的诫命的人，应受绝罚。因为主论及诫命的果实曾说：\BibleQuote{没有我，你们什么也不能做}，而不是\Quote{没有我，你们能做，只是困难}。\par

% structure: p0508 source=h3 target=subsection reason=relative-heading-rank; base=h2

\subsection{第一一四条（希腊编号一一五）}

\Emph{圣人们所说的\Quote{我们说没有罪，就欺骗了自己}这句话，不仅是谦卑的话，也是真实的话}\par

同样，议定：正如圣若望宗徒所说：\BibleQuote{如果我们说我们没有罪，就是欺骗自己，真理就不在我们内}，凡认为这句话应理解为出于谦卑我们才该说有罪，而不是因为事实如此的人，应受绝罚。因为宗徒接着又说：\BibleQuote{如果我们承认我们的罪，祂是忠信和公义的，必赦免我们的罪，并洗净我们的一切不义}，由此足以清楚看出，这话不仅是出于谦卑，也是真实的。因为宗徒本可以说：\Quote{如果我们说我们没有罪，就会自高自大，在我们内就没有谦卑}；但当他说：\Quote{我们欺骗自己，真理就不在我们内}时，便充分表明，那说自己没有罪的人所说的不是真话，而是假话。\par

% structure: p0511 source=h3 target=subsection reason=relative-heading-rank; base=h2

\subsection{第一一五条（希腊编号一一六）}

\Emph{圣人们在主祷文中为自己说：\Quote{求你宽恕我们的过犯}}\par

议定：凡说圣人们在主祷文中说\BibleQuote{求你宽恕我们的过犯}时，不是为自己说的（因为他们不需要这个祈求），而是为其余有罪的子民说的；因此，没有一个圣人可以说\Quote{求你宽恕我的过犯}，而只能说\Quote{求你宽恕我们的过犯}；这样，义人被认为是为此别人而不是为自己祈求的人，应受绝罚。因为圣雅各伯宗徒是圣洁而公义的，他说：\BibleQuote{因为我们在许多事上全都得罪了天主}。为什么加上了\Quote{全都}，除非是让这句话也符合圣咏中所读到的：\BibleQuote{上主，求你不要审问你的仆人，因为在祢面前没有一个活人是正义的}；以及极智慧的撒罗满的祈祷：\BibleQuote{没有不犯罪的人}；还有圣约伯书：\BibleQuote{祂在每人的手上封印，为使人人都能认识自己的软弱}；因此，连圣洁而公义的达尼尔在祈祷时也多次说：\BibleQuote{我们犯了罪，我们行了不义}，以及其他他真切而谦卑地承认的话；也不要以为（如有些人所想的）这话不是为自己的罪，而是为百姓的罪说的，因为他后来又说：\BibleQuote{当我祈祷，承认我的罪和我百姓的罪于上主我的天主时}；他不愿说\Quote{我们的罪}，而说\Quote{我百姓的罪和我自己的罪}，因为他作为先知，预见后来的人会这样误解他的话。\par

% structure: p0514 source=h3 target=subsection reason=relative-heading-rank; base=h2

\subsection{第116条。（希腊编号cxvii。）}

\Emph{圣人们准确地说：\Quote{宽恕我们的罪过}}\par

同样，会议也决定：任何人若认为当我们念主祷文中的这些话——\Quote{宽恕我们的罪过}——时，圣人们是出于谦卑而非出于真实而说的，就让他被弃绝。因为谁会向天主而非向人做说谎的祈祷呢？谁会用嘴唇说他希望自己的罪得到赦免，心里却认为自己没有罪需要赦免呢？\par

% structure: p0517 source=h3 target=subsection reason=relative-heading-rank; base=h2

\subsection{第117条。（希腊编号cxviii。）}

\Emph{关于从\OriginalTerm{多纳徒派}{Donatists}皈依的人民}\par

同样，会议也决定：由于几年前一次全体会议已如此规定，凡在针对多纳徒派的法律颁布之前、当这些教会成为公教会时可被建立的任何教会，都应属于那些促使其回归公教合一的主教的教座；但在此法律之后，任何曾共融的教会都应属于它们仍为多纳徒派时所归属的地方。然而，由于后来主教之间就教区发生了许多争议，且现在仍在不断涌现，而这些在当时完全无法预见，本会议现在决定：无论何处有分属不同教座的公教派和多纳徒派，无论合一是在法律颁布之前还是之后达成，教会都应属于原有公教会所在的教座。\par

% structure: p0520 source=h3 target=subsection reason=relative-heading-rank; base=h2

\subsection{第118条。（希腊编号cxix。）}

\Emph{公教会主教与从\OriginalTerm{多纳徒派}{Donatists}皈依的主教应如何划分\OriginalTerm{教区}{dioceses}}\par

同样，会议也决定：若有一位从多纳徒派皈依公教合一的主教，则应在原先双方并存的地方进行均分；即一些地方归一方，另一些归另一方；分割应由任职主教职时间较长者进行，并由较年轻者选择。但若只有一个地方，则归距离较近的一方。若两地距离与两个主教座堂相等，则归人民所选择的一方。若老公教徒希望有自己的主教，而皈依的多纳徒派也如此，则以人数较多一方的意愿为准；若双方人数相等，则归任职主教时间最长的。若发现双方并存的地方很多，以致无法均分，例如数目不等，则在分配了等数的地方后，剩余的一处应按上述只有一处地方的情况处理。\par

% structure: p0523 source=h3 target=subsection reason=relative-heading-rank; base=h2

\subsection{第119条。（希腊编号cxx。）}

\Emph{若主教占有他从\OriginalTerm{异端}{heresy}夺得的教区三年，则无人可从他手中取走}\par

同样，会议也决定：若有人在法律颁布之后将任何地方引至公教合一，并在三年内未受反对地保有之，此后不得从他手中夺走。然而，若在这三年内有一位主教可以要求该地却保持沉默，他将失去机会。但若没有主教，则不对教座造成损害；但若该地当时无主教，此后获得主教，则可自该日起三年内提出要求。同样，若一位多纳徒派主教皈依公教一方，则过去的时光不计入对他不利，而应从其皈依日起三年内，他有权利对其教座所属的地方提出要求。\par

% structure: p0526 source=h3 target=subsection reason=relative-heading-rank; base=h2

\subsection{第120条。（希腊编号cxxi。）}

\Emph{关于那些在无持有人同意的情况下，侵入他们认为属于自己的人民的人}\par

同样，会议也决定：凡主教追求他们视为属于自己教座的人民，不是通过将案件提交主教法官，而是趁他人占据该地时强行闯入者，无论这些人民是否愿意接受他们，他们都将丧失案件。凡如此行事者，若两主教之间的争端尚未完结而仍在进行，则让未经教会法官判决而侵入者离开；任何人不得以从首席主教处获得信函为理由而沾沾自喜以为能保留所夺之物；不论他有没有这样的信函，持有和接受信函者应当证明确实和平地拥有了属于他的教会。但若他提交了任何问题，则应由主教法官裁决，无论这些法官是首席主教为他们指定的，还是他们双方同意选出的邻近主教。\par

% structure: p0529 source=h3 target=subsection reason=relative-heading-rank; base=h2

\subsection{第121条。（希腊编号cxxii。）}

\Emph{关于那些忽视属于自己的人民的人}\par

同样，会议也决定：凡忽视将其教座所属的地方引至公教合一者，应由邻近的勤勉主教告诫他们不可再拖延；若自告诫之日起六个月内仍无所作为，则这些地方应归能赢得它们的人所有；但附加以下条件：若被认为自然拥有这些地方的人能够证明这种忽视是故意的，且比他人表面的更大勤奋更为有效，则当主教法官确信如此时，应将这些地方归还其教座。若争端双方是不同教省的主教，则应由该地所属教省的首席主教授命法官；但若双方同意选择邻近主教为法官，则可选择一或三位：若选择三位，则可遵循所有或两位的判决。\par

% structure: p0532 source=h3 target=subsection reason=relative-heading-rank; base=h2

\subsection{第122条。（希腊编号cxxiii。）}

\Emph{所选法官的判决不应被轻视}\par

对于双方同意所选出的法官，不得上诉；凡被发现提起此类上诉且固执地不愿服从法官者，当此事向首席主教证实后，首席主教应发出信函，使任何主教不得与他共融，直到他屈服为止。\par

% structure: p0535 source=h3 target=subsection reason=relative-heading-rank; base=h2

\subsection{第123条。（希腊编号cxxiv。）}

\Emph{若主教疏忽其教区，应剥夺其共融}\par

若在母座堂中，一位主教对异端疏忽，则应召集邻近的勤勉主教开会，向他指出其疏忽，使其无可推诿。但若在会议后六个月内，若其本省有执行者，而他未设法使他们归向大公教会合一，则无人可与祂共融，直到他尽其职责。但若无执行者来到该处，则不应归咎于该主教。\par

% structure: p0538 source=h3 target=subsection reason=relative-heading-rank; base=h2

\subsection{第124条。（希腊第125条）}

\Emph{论主教在有关多纳徒派共融事宜上说谎}\par

若经证实某位主教就那些人（原属多纳徒派）的共融说谎，并曾说他们已共融，而祂明知他们并未共融，则祂应丧失其主教职。\par

% structure: p0541 source=h3 target=subsection reason=relative-heading-rank; base=h2

\subsection{第125条。（希腊第126条）}

\Emph{论司铎及圣职人员除向非洲主教会议外不得上诉}\par

此外，会议决定：应受审判的司铎、执事或其他下级圣职人员，若质疑其主教的决定，则经其主教同意，由他们邀请邻近主教审理，并裁决其间一切争端。但若他们认为应提出上诉，则不得向海外提出上诉，仅可向非洲会议或本省首席主教提出。凡试图向海外上诉者，在非洲不得与任何人共融。\par

% structure: p0544 source=h3 target=subsection reason=relative-heading-rank; base=h2

\subsection{第126条。（希腊第127条）}

\Emph{论童贞女，即使未成年，也应给予覆面礼}\par

此外，会议决定：无论哪位主教，因童贞纯洁的危险所必需，即当惧怕有权势的求婚者或某种强夺者，或她因某些死亡顾虑而可能未覆面死去时，经其父母或受委托照顾她的人请求，可对一位童贞女给予覆面礼，或在她未满二十五岁时已给予覆面礼，规定该年龄的会议不得反对他。\par

% structure: p0547 source=h3 target=subsection reason=relative-heading-rank; base=h2

\subsection{第127条。（希腊第128条）}

\Emph{论主教在会议中不被耽搁过久，他们应从各自省份中选出三位审判员}\par

此外，会议决定：为免所有出席主教被长久耽搁，全体会议应从各省选出三位审判员。他们为迦太基省选出文森特、福尔图纳提安和克拉鲁斯；为努米底亚省选出阿利比乌斯、奥古斯丁和雷斯提图图斯；为比扎塞纳省选出圣善长者多纳提安首席主教、克雷斯孔尼乌斯、约昆都斯和埃米利安；为毛里塔尼亚·西特芬西斯省选出塞维里安、阿西亚提库斯和多纳图斯；为的黎波里塔尼亚省选出普劳提乌斯，他依惯例单独作为特使被派遣。所有这些人与圣善长者奥勒留共同审理一切事宜，全体会议请求奥勒留签署会议所做的一切，无论法令或信函。他们签名如下：我，迦太基教会主教奥勒留，同意此法令，阅读后签名。同样，全体也签名。\par

这里提到两位圣善长者，我们确信两人都是首席主教。参看第100条（第104条）。\par

\Emph{此外，本次会议有罗马教会使团出席}\par

最荣耀的皇帝霍诺留第十二次执政，狄奥多西第八次执政，奥古斯都之后，六月朔日，在迦太基，修复的巴西利卡会议室中，当奥勒留主教与罗马教会特使、意大利皮塞农省波滕提亚教会主教福斯提努斯，总督省特使文森特·卡尔维塔（库洛西塔努斯）、那不勒斯主教福尔图纳提安、马里安努斯·乌齐帕伦西斯、西米迪卡主教阿德奥达图斯、卡尔皮主教彭塔迪乌斯、穆祖巴主教鲁菲尼安、西西里主教普雷特克斯塔图斯、维里（维伦西斯）主教奎奥德武尔特代乌斯、阿比里塔主教坎迪杜斯、乌提卡主教加洛尼安；努米底亚省特使塔加斯特主教阿利比乌斯、希波勒吉亚主教奥古斯丁和卡拉马主教波西多尼乌斯；比扎塞纳省特使阿奎主教马克西米安、苏费图拉主教约昆都斯和霍雷亚-卡西利亚主教希拉里；毛里塔尼亚·西提芬西斯省特使西提菲主教诺瓦图斯和莫克塔主教利奥；毛里塔尼亚·凯撒利恩西斯省特使鲁苏卡鲁姆主教尼内卢斯、伊科西乌姆主教劳伦斯和鲁斯古尼乌姆主教努梅里安；全体会议选出的审判员就座，执事侍立，并在完成某些事项后，许多主教抱怨无法等待其余待办事务完成，且须赶回各自教会；全体会议决定，各省应选出一些人留下，以完成剩余事务。于是，那些签名证实出席者留了下来。\par

% structure: p0553 source=h3 target=subsection reason=relative-heading-rank; base=h2

\subsection{第128条。（希腊第129条）}

\Emph{论那些不共融者不应被允许提出控告}\par

全体决定，正如先前会议所规定，关于哪些人可被接纳对圣职人员提出控告；由于此前未明确哪些人不应被接纳，故我们定义：凡已被绝罚且仍处于该谴责之下者，无论欲作控告者是圣职人员或平信徒，均不可被适当接纳提出控告。\par

419年迦太基会议在其5月25日的首次会议中已完成上述内容。但当它在同月30日再次开会时，继续了法规。关于这次新会议的介绍就是本介绍。随后制定的法令是原本的，即第128、129、130、131、132和133条。\par

% structure: p0557 source=h3 target=subsection reason=relative-heading-rank; base=h2

\subsection{第129条。（希腊第130条）}

\Emph{论奴隶、释奴及一切名誉不良者不应提出控告}\par

众人皆以为宜：凡奴隶或确切意义上的自由人，均不得被接纳为控告者；凡依公法禁止在刑事诉讼中提出控告者，亦同。此亦适用于所有带有耻辱污点的人，即演员、行卑劣之事者，以及异端者、外教人或犹太人；但所有被剥夺控告权的人，并不禁止他们为自己的案件提出控告。\par

% structure: p0560 source=h3 target=subsection reason=relative-heading-rank; base=h2

\subsection{第一三〇条（希腊第一三一条）}

\Emph{未能证明一项指控者，不得被允许就另一项指控作证}\par

同样，众人以为宜：每当指控者向圣职人员提出多项罪行，且其中最先审查的一项无法证明时，不得允许他们继续就其他指控作证。\par

% structure: p0563 source=h3 target=subsection reason=relative-heading-rank; base=h2

\subsection{第一三一条（希腊第一三二条）}

\Emph{应允许谁作证}\par

凡被禁止作为控告者的人，也不得作为证人出庭；亦不得接纳指控者从其家中带来的人。凡未满十四岁者，不得被接纳作证。\par

% structure: p0566 source=h3 target=subsection reason=relative-heading-rank; base=h2

\subsection{第一三二条（希腊第一三三条）}

\Emph{关于主教排除一位自称仅向主教告明其罪行的人于共融之外}\par

亦以为宜：若某主教有时声称某人曾单独向其告明某项个人罪行，而此人现在否认，则主教不得因仅凭其个人言语未被相信而认为受轻视，尽管他声称出于良心的顾虑，不愿与此否认者共融。\par

% structure: p0569 source=h3 target=subsection reason=relative-heading-rank; base=h2

\subsection{第一三三条（希腊第一三四条）}

\Emph{主教不应轻率剥夺任何人的共融}\par

只要某人的本主教不与受绝罚者共融，其他主教也不应与该主教共融，以便主教更加谨慎，不指控他人无法向其他人以书面证据证明之事。\par

% structure: p0572 source=h3 target=subsection reason=relative-heading-rank; base=h2

\subsection{（希腊第一三五条）}

奥勒留主教说：依照本次全体公会议的章程，以及鄙人的意见，以为宜结束前述全部标题的所有事项，并让教会文件接受今日议定之讨论。\par

凡尚未表达（希腊文作\Quote{论及}）之事，我们将在次日通过我们的弟兄福斯提诺主教及长老斐理伯与阿塞卢斯，致函我们可敬的弟兄及同僚博尼法斯主教；他们以书面表示同意。\par

% structure: p0575 source=h3 target=subsection reason=relative-heading-rank; base=h2

\subsection{第一三四条（希腊文第一三五条的延续）}

\Emph{以下是自全非洲公会议致罗马城主教博尼法斯的信函，由主教福斯提诺以及长老斐理伯与阿塞卢斯（罗马教会代表）转达}\par

致最可敬的主、我们可敬的弟兄博尼法斯：奥勒留、努米底亚首席主教座瓦伦廷，以及与我们在全非洲公会议中的其他二一七位与会者。\par

既然主乐意让我们的卑微写信，论及我们神圣的弟兄、同为主教的\Person{福斯丁努}以及同为司铎的\Person{斐理伯}和\Person{亚塞鲁}与我们共行之事，并非致予有福纪念的\Person{佐西姆}主教——他们从该主教处带来命令和信件给我们——而是致予你的圣德，你由天主权威在他位置上被设立；我们应当简要陈述双方同意所决定之事；并非那些载于冗长卷宗中的事，在其中虽保持了爱德，但我们并非没有少许争论之劳顿，而拖延了商议那些现在与案件相关的事。然而，我主与弟兄，若他仍在肉身内，他将更感激地接受这消息，如同他将见到此事更和平的结局！司铎\Person{阿皮亚里}，关于其祝圣、绝罚和上诉，不仅在\Place{锡卡}，而且在全非洲教会引起了不小的丑闻，在他为所有罪过寻求宽恕后，已恢复共融。首先，我们的同僚\Place{锡卡}主教\Person{乌尔班}无疑纠正了他身上似乎需要纠正的一切。因为应牢记教会现在和未来的平安与宁静，既然此前已有如此多类似的恶事，就必须注意防止将来发生类似甚至更严重的恶事。我们认为，司铎\Person{阿皮亚里}应从\Place{锡卡}教会调离，只保留其等级的尊荣，并可在任何他愿意且能够的地方行使司铎的职务，并为此获得信件。这我们在他信中提出的请求下毫不困难地允准了。但在案件如此了结之前，在我们每日讨论处理的其他事中，案件的性质要求我们请求我们的弟兄、同为主教的\Person{福斯丁努}以及同为司铎的\Person{斐理伯}和\Person{亚塞鲁}，说明他们被委托与我们商议的事，以便纳入教会记录。他们开始口头陈述，但我们恳切要求他们以书面提出时，他们便出示了\OriginalTerm{备忘录}{Commonitory}。这份备忘录被宣读给我们，并载入记录中，他们将把这些记录带给你们。在这备忘录中，他们受命与我们商议四件事：第一，关于主教向罗马教会教宗的申诉；第二，主教不应不合宜地航海至宫廷；第三，关于若司铎和执事被其主教错误绝罚，应由邻近主教审理其案件；第四，关于\Person{乌尔班}主教，他应被绝罚甚至送至\Place{罗马}，除非他纠正了似乎需要纠正之处。关于所有这些事，关于第一和第三点，即允许主教向\Place{罗马}申诉，以及神职人员的案件应由本省主教解决，去年我们已费心在致同一可敬纪念的\Person{佐西姆}主教的书信中暗示，我们愿意暂时遵守这些规定，不伤害他，直到尼西亚大公会议的法令查寻完毕。现在我们恳求你的圣德，使你令我们遵守我们教父在尼西亚大公会议上的行为和法令，并令你那里执行他们备忘录中所带来之事：即，若某主教被控告，等等。［此处应插入\Emph{撒迪卡第七法令}］\par

同样论及司铎与执事。若某主教轻易动怒，等等。［此处应插入\Emph{撒迪卡第十七法令}］\par

这些事项已载入会议记录，直至最准确的尼西亚大公会议文本送达。若这些规定确载于其中（正如我们弟兄们从宗座寄给我们的备忘录所声称的），并且你们在意大利也依此秩序遵行，那么我们绝不可能被迫忍受我们不愿提及的对待，也无法承受不堪忍受之事。但我们相信，藉着我们的主天主的仁慈，在你圣座主持罗马教会期间，我们将不必遭受那种骄傲（\OriginalTerm{那种骄傲}{istum typhum passuri}）。并且，对我们这些不作争论的人，将保持兄弟之爱所应持守的态度。你们也将按照至高者所赐的智慧与正义，明察当如何遵守——倘若尼西亚大公会议的法典（与你们所认为的）有所不同。因为我们虽阅读过极多抄本，但在拉丁抄本中从未见过前述备忘录所载的此类谕令。同样，由于我们在此也无法在任何希腊文本中找到它们，我们曾希望从东方教会取得法典抄本，因为据说那里可找到正确的抄本。为此，我们恳请尊驾也惠予写信给那些地区的教长，即安提约基雅、亚历山大里亚和君士坦丁堡的教会，以及其他任何教会（若蒙尊驾乐意），以便那些由教父们在尼西亚城制定的同一法典能送达我们，从而藉主的帮助，为西方所有教会带来这一极大益处。因为谁能怀疑在希腊帝国收集的尼西亚大公会议抄本最为准确呢？这些抄本出自如此众多且如此尊贵的希腊教会，相互比较时竟完全一致！在此事完成之前，我们愿意允许遵守前述备忘录中为我们所定的规定——关于主教向罗马教会教长的上诉，以及关于圣职人员的案件应由其本省主教了结——直至证明送达，并相信你的真福将按天主旨意在此事上帮助我们。至于我们主教会议中所处理及定义的其他事项，由于前述弟兄、我们的同僚主教福斯蒂诺以及司铎斐理伯和阿塞鲁斯正携带会议记录同行，若你惠予接收，将使你的圣座知晓。他们签署了。\Signature{阿利比乌、奥斯定、波西狄乌、马里努斯以及其余二百一十七位主教也签署了。}\par

% structure: p0581 source=h3 target=subsection reason=relative-heading-rank; base=h2

\subsection{第一三五条法令（希腊文本中未编号）}

\Emph{以下是亚历山大里亚主教济利禄致非洲主教会议的回函，其中他发送了尼西亚大公会议的真实记录，由司铎依诺森从希腊文译出。这些信函连同同一尼西亚大公会议的文件，于四一九年十二月第六日前一天，通过前述司铎依诺森及迦太基教会副辅祭玛尔塞鲁，送至罗马教会主教圣波尼法爵。}\par

致最可敬的主公、我们圣洁的弟兄及同僚主教奥勒留、瓦伦提努，以及全体在迦太基集会的圣洁主教会议：济利禄在主内问候你们的圣座。\par

我满怀喜乐地由我儿、司铎依诺森手中，接获了尊驾如此虔诚的信函。信中你们希望，我们将从本教会档案中，把在比提尼亚都会尼西亚举行的圣父们的主教会议法令最准确的抄本，连同我们自己确认准确的证明书发送给你们。作为对此请求的回应，最可敬的主公及弟兄们，我认为有必要通过携此信的我儿、司铎依诺森，向你们发送在比提尼亚尼西亚召开的主教会议决议的最忠实抄本，并致问候。当你们查阅教会历史时，也会在那里找到这些文件。关于复活节，如你们所写，我们通知你们：我们将在下一个小纪五月第十八日之前庆祝。署名。\Signature{愿天主及我们的主照我们所愿，保存你们的圣洁主教会议，亲爱的弟兄们。}\par

% structure: p0585 source=h3 target=subsection reason=relative-heading-rank; base=h2

\subsection{第一三六条法令（希腊文本中未编号，但有新标题）}

\Emph{以下是君士坦丁堡主教阿提库致同一会议的书信开始}\par

致我们圣洁的主公、理当最蒙真福的弟兄及同僚主教奥勒留、瓦伦提努，以及在迦太基召开的主教会议中集会的其他诸位爱者：主教阿提库。\par

通过我儿、副辅祭玛尔塞鲁，我满怀感谢接获了尊驾的文书，赞美主，我得以享受如此众多弟兄的祝福。我主及最蒙真福的弟兄们，你们写信请我寄送最准确的、教父们在比提尼亚都会尼西亚城为阐明信仰所制定的法令抄本。有谁会拒绝向弟兄们提供共同的信仰或教父们所制定的法规呢？因此，我通过匆匆赶路的同一位我儿，即你们的副辅祭玛尔塞鲁，向你们寄送教父们在尼西亚城所采纳的全部法令；并请求你们，你们圣洁的主教会议常在祈祷中纪念我。署名。\Signature{愿我们的天主照我们所愿，保存你们的圣德，至圣的弟兄们。}\par

% structure: p0589 source=h3 target=subsection reason=relative-heading-rank; base=h2

\subsection{第一三七条法令（在希腊文本中是前一条的续篇）}

\Emph{以下是尼西亚大公会议的范本，于四一九年十二月第六日前一天——最荣耀的皇帝奥诺留第十二次任执政官、德奥多西第九次任执政官之后——发送给罗马城主教波尼法爵。}\par

\Quote{我们信唯一的天主……公教会与宗徒教会咒逐他们。}\par

在此信仰标志之后，附有来自前述教长的同一尼西亚大公会议法令的抄本，其内容与以上所载完全一致。我们认为，此处无需重新写出。\par

% structure: p0593 source=h3 target=subsection reason=relative-heading-rank; base=h2

\subsection{第一三八条法令（希腊文本中未编号）}

\Emph{以下是非洲主教会议致罗马城主教则肋斯定的书信开始}\par

致主、至爱及我们可敬的弟兄\Person{塞莱斯廷}、\Person{奥勒留}、\Person{帕拉蒂努斯}、\Person{安东尼}、\Person{托图斯}、\Person{塞尔乌斯代}、\Person{特伦提乌斯}、\Person{福尔图纳图斯}、\Person{马丁}、\Person{雅努阿里乌斯}、\Person{奥普塔图斯}、\Person{塞蒂西乌斯}、\Person{多纳图斯}、\Person{特阿西乌斯}、\Person{文森特}、\Person{福尔图纳提安}，以及其他聚集在\Place{迦太基}非洲全体会议的同仁们。\par

我们多么希望，正如尊座在由我们的司铎同仁\Person{利奥}送来的信中所告知我们的，您对\Person{阿皮亚里乌斯}的到来感到欣喜，同样我们也能愉快地将他澄清的书信呈递给您。那么，真实地，我们自己的满意和您后来的满意都会更加合理；而您最近表达的关于他未来听审以及已过去的听审，就不会显得仓促和欠考虑了。于是，在我们的圣弟兄和同仁主教\Person{福斯蒂努斯}到来之际，我们召集了一个会议，并相信他是与那个人一同被派来，以便正如他[阿皮亚里乌斯]先前在他的协助下被恢复司铎职，现在也可以在他的努力下，从\Place{塔布拉卡}居民指控他的极大罪行中得到澄清。但是，我们会议中正当的审查过程发现了他如此重大且骇人的罪行，以至于压倒了\Person{福斯蒂努斯}，他行事更像是一名前述之人的辩护者而非法官，并且战胜了更多是辩护者的热情而非调查者的公正。因为他首先激烈地反对整个集会，对我们施加了许多伤害，借口维护罗马教会的特权，并希望他被我们接纳到共融中，理由是尊座相信他已上诉（尽管无法证明），从而恢复了他的共融。但我们绝不允许，您通过阅读会议记录也会更好地看到。然而，经过一场持续三天的极其辛劳的调查，在此期间我们带着极大的痛苦审理了针对他的各种指控，天主——公义、刚强且忍耐的审判者——突然一击，终止了我们同仁主教\Person{福斯蒂努斯}的拖延和\Person{阿皮亚里乌斯}本人的推诿，他企图借此掩盖其可憎的恶行。因为他顽强而不知羞耻的固执被克服了，他企图通过无耻的否认来掩盖他欲望的污泥，天主如此作用于他的良心，并且公开了——甚至在人眼前——那些他已经在内心定罪的秘密罪行（一个罪恶的猪圈），以致在他虚假否认之后，他突然爆发，供认了所有他被指控的罪行，最终自愿地定了自己所有难以置信的丑行之罪，并将我们曾怀有的希望转变为呻吟，我们曾相信并希望他得以从这些可耻的污点中澄清，除了它确实在某种程度上缓解了我们的悲伤，因为他使我们免于更长时间调查的辛劳，并通过告解给他的伤口涂抹了某种补救，但主内兄弟，这是不情愿的，并且是带着挣扎的良心做出的。因此，先向您致以应有的问候，我们恳切地恳求您，今后不要轻易接纳从我们这里来的人加以听审，也不要选择将那些被我们绝罚过的人收纳到您的共融中，因为可敬的阁下，您将很容易看出，这也是尼西亚会议所规定的。因为尽管这似乎在那里针对下级圣职人员或平信徒被禁止，更何况它愿意在主教的情况下也遵守此规定，以免那些在自己的教省中被暂停共融的人，可能被尊座仓促或不恰当地恢复共融。愿尊座如您所配得的，拒绝同样司铎以及下级圣职人员那种无原则地向您寻求庇护的行为，因为教父们没有任何法令剥夺了非洲教会这一权力，而且尼西亚法令最明确地将不仅下级圣职人员，而且主教本身，都托付给了他们自己的都主教。因为他们以大智慧与公正规定，所有事务都应在发生之处终结；他们并不认为圣神的恩宠会缺少给任何一个教省，以使基督的\OriginalTerm{主教}{Sacerdotibus}明智地辨别并坚定地维护正义；尤其因为任何人若认为自己在任何判决中受冤，可以向他的教省议会上诉，甚至向全体会议\EditorialNote{即非洲会议}上诉，除非有人想象天主能将正义启迪给单独一人，而拒绝给无数\OriginalTerm{主教}{sacerdotum}的议会集会。我们怎能依赖海外做出的判决呢？因为无法将必要的证人送到那里，不论是由于性别的软弱，或年事已高，或任何其他障碍。至于尊座应从您方面派遣任何人，我们找不到任何教父会议的规定。因为关于您通过我们同一位主教弟兄\Person{福斯蒂努斯}送来的、据称包含在尼西亚会议中的内容，我们在该会议更可靠的抄本中找不到任何此类内容；这些抄本是由我们的弟兄、亚历山大里亚教会的主教圣\Person{西里尔}，以及可敬的君士坦丁堡主教\Person{阿提库斯}送给我们的，我们先前曾通过司铎\Person{依诺森}和副助祭\Person{玛尔谷}（我们通过他们收到了这些抄本）送给了您的可敬前任\Person{博尼法斯}主教。此外，无论谁希望您委派您的任何圣职人员来执行您的命令，都不要答应，以免显得我们正在把世俗统治的骄傲引入基督的教会，这教会向所有渴望看见天主的人展示出纯朴的光明和谦卑的白天。因为既然可怜的\Person{阿皮亚里乌斯}已因他可怕的罪行被逐出基督的教会，我们对于我们的弟兄\Person{福斯蒂努斯}充满信心，相信通过尊座的公正与温和，非洲将不必再忍受他，而不违反兄弟之爱。主内兄弟，愿我们的主长久保佑尊座为我们祈祷。\par

\SourceRecord{Translated by Henry Percival.；Nicene and Post-Nicene Fathers, Second Series；Vol. 14.；Edited by Philip Schaff and Henry Wace.；Buffalo, NY: Christian Literature Publishing Co.,；1900.；<http://www.newadvent.org/fathers/3816.htm>.}
